\documentclass[a4paper,12pt,oneside]{scrbook}
\KOMAoptions{
    chapterprefix=true,
    parskip=half,
    captions=tableheading,
    draft=false
}

\usepackage{iftex}

% Activar números de línea
% \usepackage{lineno}
% \linenumbers
% \setlength\linenumbersep{5pt}
% \renewcommand\linenumberfont{\normalfont\tiny\sffamily\color{gray}}

\usepackage[spanish,es-nolists,es-tabla,es-noindentfirst]{babel}
\usepackage[hidelinks]{hyperref}

\renewcommand{\listtablename}{Índice de tablas}    
\renewcommand{\listfigurename}{Índice de figuras}
\newcommand*{\listofloaname}{Índice de algoritmos}

\addto\extrasspanish{
    \renewcommand{\sectionautorefname}{Apartado}
    \renewcommand{\subsectionautorefname}{Apartado}
    \renewcommand{\subsubsectionautorefname}{Apartado}
}

% Márgenes del documento
\usepackage{geometry}
\geometry{left=2cm, right=2cm, top=2cm, bottom=3.3cm}

% Interlineado 1.5
% \onehalfspacing

% Sangría 
% Ojo, la guía de KOMA-Script sugiere usar espaciado o sangría para delimitar los párrafos,
% pero no ambas cosas porque es redundante. Descomentando la siguiente línea se activa la sangría.
% \setlength{\parindent}{0.5cm}

% Configuración de fuentes
\ifPDFTeX
    \usepackage[T1]{fontenc}
    \usepackage[utf8]{inputenc}
    \usepackage{lmodern}
    \usepackage{amsmath}
    \usepackage{amssymb}
    \usepackage{textcomp}
    \usepackage[scale=0.85]{sourcecodepro}

    \renewcommand{\familydefault}{\sfdefault}
    \newcommand{\lmromanfont}{\rmfamily}
\else % XeLaTeX / LuaLaTeX
    \usepackage{fontspec}
    \usepackage{amsmath}
    \usepackage{unicode-math}   % Incluye Latin Modern Math
    
    \setmainfont{Latin Modern Sans}[
        Ligatures=TeX,
        % Necesario porque Latin Modern Sans no incluye estilo small caps
        SmallCapsFont=Latin Modern Roman Caps
    ]
    \setsansfont{Latin Modern Sans}[Ligatures=TeX]
    \setmonofont{Source Code Pro}[Scale=MatchLowercase]
    % \setmonofont[Scale=MatchLowercase]{Fira Mono}
    \newfontfamily{\lmromanfont}{Latin Modern Roman}
\fi

% Configuración de listas (bullets — \setlist se aplica tras cargar enumitem)
\setkomafont{itemizelabel}{\lmromanfont}
\renewcommand{\labelitemi}{\textbullet}
\renewcommand{\labelitemii}{$\circ$}
\renewcommand{\labelitemiii}{$\diamond$}
\renewcommand{\labelitemiv}{$\triangleright$}

% Citas
\usepackage{csquotes}
\renewcommand{\mkbegdispquote}[2]{«}
\renewcommand{\mkenddispquote}[2]{#1»#2}

%% Bibliografía
% Numérica (ordenada por orden de aparación)
\usepackage[style=numeric-comp,sorting=none,backend=biber]{biblatex}
\defbibnote{note}{Las siguientes referencias bibliográficas se presentan en orden de aparición en el texto.\par\bigskip}
% APA (ordenada por autor, año y título)
% \usepackage[style=apa,sorting=nyt,backend=biber]{biblatex}
% \defbibnote{note}{Las siguientes referencias bibliográficas se presentan en orden alfabético por autor. Las referencias con más de un autor aparecen ordenadas en base al primero de los mismos.\par\bigskip}

\addbibresource{referencias.bib}

\usepackage{caption}
\usepackage{enumitem}
\setlist{leftmargin=1.5em}
\usepackage{float}
\usepackage{graphicx}
\usepackage{lipsum}
\usepackage{microtype}
\usepackage{pdflscape}
\usepackage{rotating}
\usepackage{subcaption}
\usepackage[table]{xcolor}

\usepackage{siunitx}
\ifPDFTeX
    \newcommand{\perthousand}{\textperthousand}
\else
    % XeLaTeX / LuaLaTeX
    \newcommand{\perthousand}{‰}
\fi

% Tablas
\usepackage{array} 
\usepackage{booktabs}
\usepackage{longtable}
\usepackage{multirow}
\usepackage{tabularx}
% Definir columnas para alineación a la derecha o centro y ajuste automático
\newcolumntype{R}{>{\raggedleft\arraybackslash}X}
\newcolumntype{C}{>{\centering\arraybackslash}X}
% Definir tipos de columnas en modo matemático
% Ver pag. 3 en la documentación del paquete array:
\newcolumntype{M}{>{$}c<{$}}
\newcolumntype{L}{>{$}l<{$}}
\newcolumntype{N}{>{$}r<{$}}
% Distancia entre filas
\renewcommand*{\arraystretch}{1.5}

% Código con coloreado de sintaxis
\usepackage{scrhack}            % Para evitar un warning sobre \float@addtolist
\usepackage{listings}
\renewcommand{\lstlistingname}{Código}
\renewcommand{\lstlistlistingname}{Índice de códigos}
\definecolor{bluekeywords}{rgb}{0.13,0.13,1}
\definecolor{greencomments}{rgb}{0,0.5,0}
\definecolor{redstrings}{rgb}{0.9,0,0}
\definecolor{punct}{rgb}{0.1,0.1,0.55}
\definecolor{delim}{rgb}{0.1,0.1,0.55}
\definecolor{numb}{rgb}{0.65,0.16,0.16}

\lstset{
  basicstyle=\ttfamily\footnotesize,  % el estilo de la fuente usado para el código
  aboveskip=12pt,                     % espacio antes del código
  belowskip=12pt,                     % espacio después del código  
  numbers=left,                       % donde poner los números de línea
  numberstyle=\tiny\color{gray},      % el estilo de los números de línea
  stepnumber=1,                       % el intervalo entre los números de línea
  numbersep=8pt,                      % cuánto espacio hay entre los números de línea y el código
  backgroundcolor=\color{white},      % el color de fondo de los cuadros de código
  showspaces=false,                   % muestra espacios agregando subrayado
  showstringspaces=false,             % subraya solo los espacios en las cadenas
  showtabs=false,                     % muestra tabulaciones agregando subrayado
  frame=leftline,                     % añade un marco alrededor del código
  framerule=0.5mm,                    % grosor de línea del marco
  rulecolor=\color{lightgray},        % si no se establece, el color del marco puede cambiar en los cuadros roto
  tabsize=2,                          % establece el tamaño de la tabulación
  captionpos=b,                       % establece la posición de la leyenda del cuadro de código
  breaklines=true,                    % líneas de ajuste automático
  breakatwhitespace=true,             % corta las líneas solo en los espacios en blanco
  keywordstyle=\color{bluekeywords}\bfseries,  % estilo para las palabras clave
  commentstyle=\color{greencomments}, % estilo para los comentarios
  stringstyle=\color{redstrings},     % estilo para las cadenas
  title=\lstname,                     % muestra el nombre de los archivos listados
  literate=
     *{0}{{{\color{numb}0}}}{1}
      {1}{{{\color{numb}1}}}{1}
      {2}{{{\color{numb}2}}}{1}
      {3}{{{\color{numb}3}}}{1}
      {4}{{{\color{numb}4}}}{1}
      {5}{{{\color{numb}5}}}{1}
      {6}{{{\color{numb}6}}}{1}
      {7}{{{\color{numb}7}}}{1}
      {8}{{{\color{numb}8}}}{1}
      {9}{{{\color{numb}9}}}{1}
      {:}{{{\color{punct}{:}}}}{1}
      {,}{{{\color{punct}{,}}}}{1}
      {\{}{{{\color{delim}{\{}}}}{1}
      {\}}{{{\color{delim}{\}}}}}{1}
      {[}{{{\color{delim}{[}}}}{1}
      {]}{{{\color{delim}{]}}}}{1},
}

\ifPDFTeX
\lstset{
  inputencoding=utf8,
  extendedchars=true,
  literate=
      {á}{{\'a}}1  {é}{{\'e}}1  {í}{{\'i}}1 {ó}{{\'o}}1  {ú}{{\'u}}1
      {Á}{{\'A}}1  {É}{{\'E}}1  {Í}{{\'I}}1 {Ó}{{\'O}}1  {Ú}{{\'U}}1
      {à}{{\`a}}1  {è}{{\`e}}1  {ì}{{\`i}}1 {ò}{{\`o}}1  {ù}{{\`u}}1
      {À}{{\`A}}1  {È}{{\`E}}1  {Ì}{{\`I}}1 {Ò}{{\`O}}1  {Ù}{{\`U}}1
      {ä}{{\"a}}1  {ë}{{\"e}}1  {ï}{{\"i}}1 {ö}{{\"o}}1  {ü}{{\"u}}1
      {Ä}{{\"A}}1  {Ë}{{\"E}}1  {Ï}{{\"I}}1 {Ö}{{\"O}}1  {Ü}{{\"U}}1
      {â}{{\^a}}1  {ê}{{\^e}}1  {î}{{\^i}}1 {ô}{{\^o}}1  {û}{{\^u}}1
      {Â}{{\^A}}1  {Ê}{{\^E}}1  {Î}{{\^I}}1 {Ô}{{\^O}}1  {Û}{{\^U}}1
      {œ}{{\oe}}1  {Œ}{{\OE}}1  {æ}{{\ae}}1 {Æ}{{\AE}}1  {ß}{{\ss}}1
      {ẞ}{{\SS}}1  {ç}{{\c{c}}}1 {Ç}{{\c{C}}}1 {ø}{{\o}}1  {Ø}{{\O}}1
      {å}{{\aa}}1  {Å}{{\AA}}1  {ã}{{\~a}}1  {õ}{{\~o}}1 {Ã}{{\~A}}1
      {Õ}{{\~O}}1  {ñ}{{\~n}}1  {Ñ}{{\~N}}1  {¿}{{?`}}1  {¡}{{!`}}1
      {„}{\quotedblbase}1 {“}{\textquotedblleft}1 {–}{$-$}1
      {°}{{\textdegree}}1 {º}{{\textordmasculine}}1 {ª}{{\textordfeminine}}1
      {£}{{\pounds}}1  {©}{{\copyright}}1  {®}{{\textregistered}}1
      {«}{{\guillemotleft}}1  {»}{{\guillemotright}}1  {Ð}{{\DH}}1  {ð}{{\dh}}1
      {Ý}{{\'Y}}1    {ý}{{\'y}}1    {Þ}{{\TH}}1    {þ}{{\th}}1    {Ă}{{\u{A}}}1
      {ă}{{\u{a}}}1  {Ą}{{\k{A}}}1  {ą}{{\k{a}}}1  {Ć}{{\'C}}1    {ć}{{\'c}}1
      {Č}{{\v{C}}}1  {č}{{\v{c}}}1  {Ď}{{\v{D}}}1  {ď}{{\v{d}}}1  {Đ}{{\DJ}}1
      {đ}{{\dj}}1    {Ė}{{\.{E}}}1  {ė}{{\.{e}}}1  {Ę}{{\k{E}}}1  {ę}{{\k{e}}}1
      {Ě}{{\v{E}}}1  {ě}{{\v{e}}}1  {Ğ}{{\u{G}}}1  {ğ}{{\u{g}}}1  {Ĩ}{{\~I}}1
      {ĩ}{{\~\i}}1   {Į}{{\k{I}}}1  {į}{{\k{i}}}1  {İ}{{\.{I}}}1  {ı}{{\i}}1
      {Ĺ}{{\'L}}1    {ĺ}{{\'l}}1    {Ľ}{{\v{L}}}1  {ľ}{{\v{l}}}1  {Ł}{{\L{}}}1
      {ł}{{\l{}}}1   {Ń}{{\'N}}1    {ń}{{\'n}}1    {Ň}{{\v{N}}}1  {ň}{{\v{n}}}1
      {Ő}{{\H{O}}}1  {ő}{{\H{o}}}1  {Ŕ}{{\'{R}}}1  {ŕ}{{\'{r}}}1  {Ř}{{\v{R}}}1
      {ř}{{\v{r}}}1  {Ś}{{\'S}}1    {ś}{{\'s}}1    {Ş}{{\c{S}}}1  {ş}{{\c{s}}}1
      {Š}{{\v{S}}}1  {š}{{\v{s}}}1  {Ť}{{\v{T}}}1  {ť}{{\v{t}}}1  {Ũ}{{\~U}}1
      {ũ}{{\~u}}1    {Ū}{{\={U}}}1  {ū}{{\={u}}}1  {Ů}{{\r{U}}}1  {ů}{{\r{u}}}1
      {Ű}{{\H{U}}}1  {ű}{{\H{u}}}1  {Ų}{{\k{U}}}1  {ų}{{\k{u}}}1  {Ź}{{\'Z}}1
      {ź}{{\'z}}1    {Ż}{{\.Z}}1    {ż}{{\.z}}1    {Ž}{{\v{Z}}}1  {ž}{{\v{z}}}1    
}
\fi

%% Rust
\lstdefinelanguage{Rust}{
  keywords=[1]{as, break, const, continue, crate, else, enum, extern, false, fn,
    for, if, impl, in, let, loop, match, mod, move, mut, pub, ref, return,
    self, Self, static, struct, super, trait, true, type, unsafe, use, where, while},
  keywords=[2]{bool, char, f32, f64, i8, i16, i32, i64, i128, isize,
    str, u8, u16, u32, u64, u128, usize, String, Vec, Option, Result,
    Box, Rc, Arc, HashMap},
  keywords=[3]{derive, cfg, allow, warn, deny, macro_rules},
  sensitive=true,
  morecomment=[l]{//},
  morecomment=[s]{/*}{*/},
  morestring=[b]",
  morestring=[b]',
}[keywords, comments, strings]

%% JavaScript
% https://github.com/ghammock/LaTeX_Listings_JavaScript_ES6
\lstdefinelanguage{JavaScript}{
  morekeywords=[1]{break, continue, delete, else, for, function, if, in,
    new, return, this, typeof, var, void, while, with},
  % Literals, primitive types, and reference types.
  morekeywords=[2]{false, null, true, boolean, number, undefined,
    Array, Boolean, Date, Math, Number, String, Object},
  % Built-ins.
  morekeywords=[3]{eval, parseInt, parseFloat, escape, unescape},
  sensitive,
  morecomment=[s]{/*}{*/},
  morecomment=[l]//,
  morecomment=[s]{/**}{*/}, % JavaDoc style comments
  morestring=[b]',
  morestring=[b]"
}[keywords, comments, strings]

\lstalias[]{ES6}[ECMAScript2015]{JavaScript}

\lstdefinelanguage[ECMAScript2015]{JavaScript}[]{JavaScript}{
  morekeywords=[1]{await, async, case, catch, class, const, default, do,
    enum, export, extends, finally, from, implements, import, instanceof,
    let, static, super, switch, throw, try},
  morestring=[b]` % Interpolation strings.
}       % Estilo y definiciones para lenguajes adicionales

% Descripción de algoritmos en pseudocódigo
\usepackage[chapter]{algorithm}
\usepackage{algpseudocodex}
\floatname{algorithm}{Algoritmo}
\captionsetup[algorithm]{labelfont=bf, labelsep=default}

\newenvironment{abstract}[1][\abstractname]{
    \cleardoublepageusingstyle{empty}
    \thispagestyle{empty}
    \begin{center}\textbf{#1}\end{center}
    \begin{itshape}\par\noindent%
}
{\end{itshape}}

\newenvironment{keywords}[1][Keywords]
{\vspace{7pt}\par\noindent\textup{\textbf{#1: }}\begin{upshape}}
{\end{upshape}}

\newcommand{\license}[2]{%
\begin{center}
    \begin{minipage}{0.8\textwidth}   
        \begin{center}
            \includegraphics[width=4.66cm]{#1}\\[12pt]
            {\Large #2}
        \end{center}
    \end{minipage}
\end{center}
}

% Formato de capítulos
% Eliminar el punto y reducir el espacio entre el prefijo y el nombre del capítulo 
\renewcommand*{\chapterformat}{%
    \mbox{\chapappifchapterprefix{\nobreakspace}\thechapter
    \IfUsePrefixLine{}{\enskip}}}
\renewcommand{\chapterheadmidvskip}{\vskip 0pt}

% Formato de los títulos en figuras y tablas
\renewcommand*{\figureformat}{\figurename~\thefigure}
\renewcommand*{\tableformat}{\tablename~\thetable}

% Redefinición de las entradas de los capítulos en la tabla de contenidos
% para que aparezca el prefijo.
\let\originaladdchaptertocentry\addchaptertocentry
\renewcommand*{\addchaptertocentry}[2]{%
  \IfArgIsEmpty{#1}{% Entrada sin número
    \originaladdchaptertocentry{#1}{#2}%
  }{% Entrada con número
    % Eliminar el número y poner el prefijo directamente en el título del capítulo
    \originaladdchaptertocentry{}{\chapapp~#1\autodot\space#2}%
  }%
}

\begin{document}   

%%%%%%%%%%%%%%%%%%%%%%%%%%%%%%%%%%%%%%%%%%%%%%%%%%%%%%%%%%%%%%%%%%%%%%%%%%%%%%%
% Portada y metadatos
%%%%%%%%%%%%%%%%%%%%%%%%%%%%%%%%%%%%%%%%%%%%%%%%%%%%%%%%%%%%%%%%%%%%%%%%%%%%%%%

% Metadatos del PDF
\hypersetup{
    pdftitle='Título del trabajo',
    pdfauthor='Nombre y Apellidos',
    pdfsubject={Trabajo de Fin de Grado},
}

\pagestyle{empty}
\newcommand{\HRule}{\rule{\linewidth}{0.3mm}}
{
    \setlength{\parindent}{0mm}
    \setlength{\parskip}{0mm}
    
    \vspace*{1.20cm}
    \includegraphics[width=9.81cm]{images/logos/escuela-ingenieria-tecnologia-original}
    \vspace*{\stretch{0.9}}
    
    {\centering
    \fontsize{32pt}{32pt}\selectfont Trabajo de Fin de Grado\\[10pt]
    \fontsize{20pt}{20pt}\selectfont Grado en Ingeniería Informática\par}
    \HRule\vspace*{-2mm}
    \begin{flushright}
        {\fontsize{32pt}{32pt}\selectfont Qalc: un lenguaje funcional para computación cuántica\par
        \vspace*{3mm}
        \fontsize{18pt}{18pt}\selectfont \textit{Qalc: a functional language for quantum computing}\par
        \vspace*{11mm}
        \fontsize{16pt}{16pt}\selectfont Guillermo Silva González
        }\vspace*{12mm}
    \end{flushright}
    \HRule
    
    \vspace*{\stretch{1}}
    \begin{center}
        \fontsize{18pt}{18pt}\selectfont La Laguna, \today
    \end{center}
}

%%%%%%%%%%%%%%%%%%%%%%%%%%%%%%%%%%%%%%%%%%%%%%%%%%%%%%%%%%
% Certificado
%%%%%%%%%%%%%%%%%%%%%%%%%%%%%%%%%%%%%%%%%%%%%%%%%%%%%%%%%%
\frontmatter
\cleardoublepageusingstyle{empty}
\thispagestyle{empty}

{
\setlength\parindent{0pt}
    D./Dña \textbf{Daniel Escanez-Exposito}, profesor/a Titular de Universidad adscrito al Departamento de Nombre del Departamento de la Universidad de La Laguna, como tutor
    
    \bigskip
    D./Dña \textbf{Pino Teresa Caballero Gil}, profesor/a Titular de Universidad adscrito al Departamento de Nombre del Departamento de la Universidad de La Laguna, como cotutor.
    
    \bigskip
    \bigskip
    \textbf{C E R T I F I C A \textcolor{red}{(N)}}

    \bigskip
    Que la presente memoria titulada:
    
    \bigskip
    \begin{quote}
    \textit{``Qalc: un lenguaje funcional para computación cuántica''}
    \end{quote}
    
    \bigskip
    \noindent ha sido realizada bajo su dirección por D./Dña \textbf{Guillermo Silva González}.
    
    \bigskip
    Y para que así conste, en cumplimiento de la legislación vigente y a los efectos
    oportunos, firman la presente en La Laguna a \today
}

%%%%%%%%%%%%%%%%%%%%%%%%%%%%%%%%%%%%%%%%%%%%%%%%%%%%%%%%%%
% Agradecimientos
%%%%%%%%%%%%%%%%%%%%%%%%%%%%%%%%%%%%%%%%%%%%%%%%%%%%%%%%%%
\cleardoublepageusingstyle{empty}
\thispagestyle{empty}

\begin{flushright}
    \setlength{\parindent}{0mm}

    {\LARGE Agradecimientos}
    \vspace{15mm}
    
    \begin{large}
        \begin{minipage}{0.3\textwidth}
            \begin{flushright}
            XXX
            
            XXX
            
            XXX
            
            XXX
            \end{flushright}
        \end{minipage}
    \end{large}
\end{flushright}

%%%%%%%%%%%%%%%%%%%%%%%%%%%%%%%%%%%%%%%%%%%%%%%%%%%%%%%%%%
% Licencia
%%%%%%%%%%%%%%%%%%%%%%%%%%%%%%%%%%%%%%%%%%%%%%%%%%%%%%%%%%
\cleardoublepageusingstyle{empty}
\thispagestyle{empty}

{\noindent\LARGE Licencia}
\vspace{15mm}

\license{images/licenses/by-sa}{© Esta obra está bajo una licencia de Creative Commons Reconocimiento-CompartirIgual 4.0 Internacional.}

%%%%%%%%%%%%%%%%%%%%%%%%%%%%%%%%%%%%%%%%%%%%%%%%%%%%%%%%%%
% Resumen
%%%%%%%%%%%%%%%%%%%%%%%%%%%%%%%%%%%%%%%%%%%%%%%%%%%%%%%%%%
\begin{abstract}
La programación cuántica actual está dominada por herramientas que operan al nivel del circuito, en las que el programador ensambla explícitamente puertas lógicas sobre cúbits individuales sin ninguna garantía formal sobre la corrección física del resultado. Bibliotecas de amplia adopción industrial como Qiskit o Cirq, al estar embebidas en un lenguaje clásico de tipado dinámico, no pueden verificar en tiempo de compilación restricciones fundamentales de la mecánica cuántica, como el teorema de no clonación o la correcta descomputación de los cúbits auxiliares; los errores derivados de su incumplimiento solo se manifiestan en tiempo de ejecución física, con un coste de diagnóstico elevado. Las propuestas académicas más rigurosas, como Silq, Guppy o Twist, incorporan alguna forma de análisis estático sobre estas restricciones, pero ninguna combina un sistema de tipos lineales completo con una salida hacia un formato estándar e interoperable con hardware de múltiples fabricantes.

Este trabajo presenta Qalc, un lenguaje de programación funcional con tipado estático fuerte orientado a la descripción de circuitos cuánticos, junto con la implementación de su compilador en Rust. El sistema de tipos de Qalc combina la inferencia de tipos de Hindley-Milner con una disciplina de tipos lineales que obliga a que cada recurso cuántico se consuma exactamente una vez a lo largo del programa, haciendo imposible en tiempo de compilación tanto la clonación accidental como la fuga de cúbits sin liberar. El modelo de cómputo del lenguaje es deliberadamente no recursivo, en correspondencia directa con la naturaleza finita y estática de los circuitos cuánticos ejecutables sobre hardware NISQ (del inglés \textit{Noisy Intermediate-Scale Quantum}), y expresa la repetición de patrones mediante combinadores de orden superior que el compilador expande en tiempo de compilación.

El compilador se estructura como un pipeline de fases independientes, desde el análisis léxico y sintáctico hasta la generación de código, pasando por la normalización, la inferencia de tipos, la verificación de recursividad, un aplanador perezoso que reduce el programa a una única expresión sin abstracción funcional pendiente, una comprobación de consumo lineal y, finalmente, la generación de código hacia OpenQASM~3 (del inglés \textit{Open Quantum Assembly Language}), el estándar abierto adoptado por los principales proveedores de hardware cuántico. El resultado de este trabajo es un prototipo funcional en el que el frontend del compilador, hasta la inferencia de tipos, está completamente implementado, mientras que las fases del backend se encuentran en un estado experimental suficiente para demostrar el funcionamiento completo del pipeline. Como líneas futuras se identifican la descomputación automática, la maduración de las fases del backend hasta cubrir la totalidad del lenguaje diseñado y la ampliación de la librería estándar de tipos y operaciones más allá del catálogo ilustrativo usado en este prototipo.

\begin{keywords}[Palabras clave]
Computación cuántica, Lenguajes de programación, Tipos lineales, Inferencia de tipos Hindley-Milner, OpenQASM, Compiladores
\end{keywords}
\end{abstract}

%%%%%%%%%%%%%%%%%%%%%%%%%%%%%%%%%%%%%%%%%%%%%%%%%%%%%%%%%
\cleardoublepageusingstyle{empty}
\thispagestyle{empty}

\begin{abstract}[Abstract]
Quantum programming today is dominated by tools that operate at the circuit level, where the programmer explicitly assembles logic gates over individual qubits with no formal guarantee about the physical correctness of the result. Widely adopted industrial libraries such as Qiskit or Cirq, being embedded in a dynamically typed classical language, cannot verify fundamental constraints of quantum mechanics at compile time, such as the no-cloning theorem or the correct uncomputation of ancilla qubits; errors arising from violating these constraints only surface at physical runtime, at a high diagnostic cost. The most rigorous academic proposals, such as Silq, Guppy or Twist, incorporate some form of static analysis over these constraints, but none combines a full linear type system with output to a standard, interoperable format supported by multiple hardware vendors.

This work presents Qalc, a statically and strongly typed functional programming language for describing quantum circuits, together with the implementation of its compiler in Rust. The type system of Qalc combines Hindley-Milner type inference with a linear type discipline that forces every quantum resource to be consumed exactly once throughout the program, making both accidental cloning and unreleased qubit leaks impossible to express at compile time. The computation model of the language is deliberately non-recursive, directly mirroring the finite, static nature of quantum circuits executable on NISQ (Noisy Intermediate-Scale Quantum) hardware, and expresses repeated patterns through higher-order combinators that the compiler unrolls at compile time.

The compiler is structured as a pipeline of independent phases, from lexical and syntax analysis to code generation, passing through normalization, type inference, recursion verification, a lazy flattener that reduces the program to a single expression free of pending functional abstraction, a linear consumption check, and finally code generation targeting OpenQASM~3 (Open Quantum Assembly Language), the open standard adopted by the main quantum hardware providers. The result of this work is a functional prototype in which the compiler frontend, up to type inference, is fully implemented, while the backend phases remain at an experimental stage sufficient to demonstrate the complete pipeline. Future work includes automatic uncomputation, maturing the backend phases to cover the full designed language, and expanding the standard library of types and operations beyond the illustrative catalog used in this prototype.

\begin{keywords}
Quantum computing, Programming languages, Linear types, Hindley-Milner type inference, OpenQASM, Compilers
\end{keywords}
\end{abstract}

%%%%%%%%%%%%%%%%%%%%%%%%%%%%%%%%%%%%%%%%%%%%%%%%%%%%%%%%%
\pagestyle{plain}
\cleardoublepage
\setcounter{page}{1} 

\tableofcontents
\listoffigures
\listoftables
\addcontentsline{toc}{chapter}{\lstlistlistingname}
\lstlistoflistings
%\listofalgorithms

%%%%%%%%%%%%%%%%%%%%%%%%%%%%%%%%%%%%%%%%%%%%%%%%%%%%%%%%%%
% Contenidos
%%%%%%%%%%%%%%%%%%%%%%%%%%%%%%%%%%%%%%%%%%%%%%%%%%%%%%%%%%
\mainmatter

\chapter{Introducción y motivación}

La computación cuántica lleva décadas prometiendo resolver problemas que escapan a la capacidad de cualquier ordenador clásico, desde la simulación de sistemas moleculares complejos hasta la búsqueda en espacios de soluciones de escala exponencial. Sin embargo, programar estas máquinas sigue siendo hoy una tarea reservada a especialistas con formación profunda en física cuántica, no porque los algoritmos sean inherentemente incomprensibles, sino porque las herramientas disponibles obligan al desarrollador a razonar al nivel del hardware. Este trabajo de fin de grado nace de esa brecha (la distancia entre lo que la computación cuántica puede hacer y lo que un ingeniero de software puede expresar con las herramientas actuales).

\section{Fundamentos de la computación cuántica}

La computación cuántica constituye un cambio de paradigma respecto a la computación clásica. Mientras que un bit clásico representa en todo momento un estado discreto (0 o 1), la unidad fundamental de información cuántica, el cúbit, puede existir en superposición, es decir, en una combinación lineal de ambos estados cuya amplitud determina la probabilidad de cada resultado al medir. Este comportamiento, descrito formalmente mediante vectores normalizados en un espacio de Hilbert complejo, permite que un sistema de $n$ cúbits represente $2^n$ estados de forma simultánea.

Sobre esta base actúan dos fenómenos adicionales que dotan a la computación cuántica de su potencia expresiva. El entrelazamiento crea correlaciones entre cúbits sin equivalente clásico, pues el estado de un cúbit entrelazado no puede describirse de forma independiente del resto del sistema. La medición, por su parte, es destructiva; al observar el estado de un cúbit, el sistema colapsa irreversiblemente hacia uno de sus valores base, destruyendo la superposición. La evolución del sistema entre mediciones se realiza mediante operadores unitarios, que conservan la norma del vector de estado y garantizan la reversibilidad de las operaciones.

Para estructurar la interacción entre el software clásico y el hardware cuántico se utiliza habitualmente el modelo QRAM (\textit{Quantum Random Access Machine}), en el que un controlador clásico envía instrucciones a un coprocesador cuántico, recibe los resultados de las mediciones y adapta el flujo de control en consecuencia. Esta arquitectura híbrida es la base sobre la que se construyen todos los lenguajes y entornos de programación cuánticos modernos.

\section{Los lenguajes de bajo nivel como barrera}

El ecosistema actual de programación cuántica está dominado por herramientas que operan al nivel del circuito cuántico, donde el desarrollador construye su algoritmo ensambladando explícitamente puertas lógicas sobre cúbits individuales. Este modelo tiene sentido como representación de bajo nivel para el hardware, pero resulta inadecuado como abstracción de programación para cualquier algoritmo de complejidad moderada.

Los entornos más extendidos, como Qiskit (IBM) o Cirq (Google), se presentan como bibliotecas Python que construyen circuitos de este modo. Su principal limitación no es la falta de funcionalidades, sino la ausencia total de garantías estáticas sobre la corrección física del código. Al ser bibliotecas de un lenguaje clásico de tipado dinámico, no pueden verificar en tiempo de compilación reglas fundamentales de la mecánica cuántica. El teorema de no clonación prohíbe duplicar el estado de un cúbit arbitrario, pero nada en Python impide que un programador asigne la misma variable cuántica a dos operaciones simultáneas; el error solo aparecerá en tiempo de ejecución, a menudo en forma de resultados silenciosamente incorrectos. La gestión manual de la topología del chip, la secuencia de pulsos o el ciclo de vida de cada cúbit traslada al programador una responsabilidad que debería pertenecer al compilador.

El resultado es un ecosistema donde la complejidad accidental supera con frecuencia a la complejidad esencial del problema. Escribir el algoritmo de Grover o la transformada cuántica de Fourier en Qiskit requiere más líneas de código de bajo nivel que de lógica algorítmica, y cualquier refactorización puede introducir errores físicos difíciles de detectar.

\section{Qalc, un lenguaje cuántico de alto nivel}

Este proyecto propone el diseño e implementación de Qalc, un lenguaje de programación cuántico de alto nivel con paradigma funcional, inspirado en la sintaxis de Haskell. El objetivo es que el desarrollador pueda expresar algoritmos cuánticos en términos de transformaciones y composiciones de datos, sin necesidad de razonar sobre registros físicos ni sobre el ensamblaje de puertas.

Qalc introduce un modelo en el que los circuitos se construyen funcionalmente; las operaciones no modifican estado global sino que devuelven nuevas representaciones del circuito en construcción. El lenguaje expone tipos de datos abstractos como \lstinline|qint| o \lstinline|qbool|, y permite componer operaciones mediante el operador pipe (\lstinline{|>}) y funciones de orden superior, de forma análoga a como un desarrollador funcional clásico encadena transformaciones sobre listas o árboles.

\subsection{Sistema de tipos fuerte con tipos lineales}

El eje de seguridad de Qalc es su sistema de tipos estático, construido sobre inferencia Hindley-Milner extendida con tipos lineales (o afines). Los tipos lineales tratan cada cúbit como un recurso físico que debe consumirse exactamente una vez en el flujo del programa: no puede copiarse, no puede ignorarse y no puede pasarse a dos operaciones a la vez.

Esta restricción tiene una consecuencia directa: el compilador puede verificar el teorema de no clonación en tiempo de compilación, antes de que el programa llegue a ejecutarse en ningún simulador ni hardware real. Cada cúbit tiene un único propietario en cada punto del flujo de datos, y cualquier violación de esa propiedad se convierte en un error de compilación con un diagnóstico preciso que señala el punto exacto del código fuente. Esto elimina una clase entera de errores que en los entornos actuales solo se manifiestan en ejecución.

\section{Alcance del proyecto}

El presente trabajo abarca el diseño completo del lenguaje y la implementación de su compilador hasta un prototipo funcional. El compilador está escrito en Rust y sigue una arquitectura de pipeline modular que comprende el análisis léxico, el análisis sintáctico con construcción del árbol sintáctico abstracto (AST, del inglés \textit{Abstract Syntax Tree}), la normalización, la inferencia de tipos con Hindley-Milner y la generación de código hacia OpenQASM~3 (del inglés \textit{Open Quantum Assembly Language}, versión 3), el estándar de descripción de circuitos cuánticos más extendido en la industria.

Quedan fuera del alcance de este prototipo, y se identifican como líneas de trabajo futuro, funcionalidades avanzadas del lenguaje como la gestión automática de la descomputación (\textit{uncomputation}). En los sistemas cuánticos, descartar un cúbit auxiliar que permanece entrelazado con el registro de cálculo induce una medición implícita que destruye la coherencia del estado global; automatizar ese proceso en el compilador, tal y como hace Silq, constituye un problema de investigación de primer orden que merece un trabajo propio.

\chapter{Diseño del lenguaje y especificación}

En este capítulo se presenta el diseño de Qalc, el lenguaje de programación cuántico de alto nivel propuesto en este trabajo. A diferencia de los capítulos de implementación, aquí se describe el lenguaje como un \textit{plano de diseño} (la especificación completa de su filosofía, su modelo semántico, su sistema de tipos, su sintaxis y su gramática formal). El objetivo es fijar la intención del lenguaje y sus reglas, con independencia de qué porción de esas reglas valida ya el prototipo del compilador. Aquellos mecanismos que forman parte del diseño pero quedan fuera del prototipo implementado se señalan explícitamente.

\section{Filosofía y objetivos de diseño}

El análisis del estado del arte realizado en el capítulo anterior evidencia una tensión recurrente en los lenguajes cuánticos actuales: los entornos de bajo nivel (Qiskit, Cirq, OpenQASM) ofrecen control físico explícito a costa de una elevada verbosidad y de la ausencia de garantías estáticas, mientras que las propuestas de alto nivel más rigurosas (Silq, Q\#, Quipper) imponen curvas de aprendizaje pronunciadas o delegan en el programador la responsabilidad de no violar restricciones físicas fundamentales. Qalc se diseña para situarse en el espacio que dejan libre estas alternativas, guiado por cuatro principios.

\paragraph{Abstracción sobre circuitos, no sobre registros físicos}
El programador no manipula registros ni índices de cúbits físicos, sino que describe transformaciones y composiciones de alto nivel. Un programa Qalc es, conceptualmente, la descripción funcional de un circuito que el compilador materializa más tarde en una representación ejecutable (OpenQASM~3 como destino por defecto). Esta decisión es el resultado de la evolución de diseño documentada durante el proyecto: el modelo inicial basado en registros físicos al estilo Qiskit se abandonó en favor de un modelo en el que las operaciones construyen el circuito de forma composicional.

\paragraph{Paradigma funcional puro e inmutable}
Qalc adopta una sintaxis y una semántica fuertemente inspiradas en Haskell: funciones de orden superior, aplicación parcial, composición mediante tuberías (\textit{pipes}) y ausencia de estado mutable global. La motivación es doble. Por un lado, los lenguajes funcionales han demostrado ser un vehículo idóneo para construir abstracciones sólidas y componibles. Por otro, la inmutabilidad encaja de forma natural con la reversibilidad de la computación cuántica unitaria.

\paragraph{Seguridad estática frente a violaciones físicas}
El sistema de tipos no es un mero mecanismo de documentación, sino la primera línea de defensa frente a programas físicamente imposibles. El objetivo de diseño es que la clase de errores más peligrosa de la programación cuántica (la clonación implícita de información cuántica, prohibida por el teorema de no clonación) sea detectable en tiempo de compilación y no en tiempo de ejecución sobre la QPU.

\paragraph{Expresividad y ergonomía}
El lenguaje busca que algoritmos no triviales (por ejemplo, la construcción de un oráculo para el algoritmo de Grover) puedan expresarse de forma concisa y legible, acercando la escritura de algoritmos cuánticos a la experiencia de la programación funcional clásica.

La tensión central que estos principios deben resolver es la aparente incompatibilidad entre la inmutabilidad funcional y la naturaleza de los circuitos cuánticos, que transforman el estado. La sección siguiente describe el modelo conceptual que reconcilia ambas.

\subsection{Modelo de cómputo no recursivo}

Una restricción de diseño fundamental, impuesta por las limitaciones del hardware cuántico actual y por el modelo de ejecución de OpenQASM~3 (el formato de destino del compilador), es que Qalc no admite recursividad. Un circuito cuántico es, en última instancia, una estructura finita y estática de puertas que debe quedar completamente determinada en tiempo de compilación, pues no existe sobre la QPU una pila de llamadas ni un mecanismo de invocación dinámica que permita que una función se invoque a sí misma un número de veces no acotado a priori. Por ello, OpenQASM~3 carece de recursión y exige que toda repetición sea acotada y desplegable estáticamente en una secuencia concreta de operaciones.

En consecuencia, Qalc es un lenguaje funcional en el que el resultado se describe mediante la composición de transformaciones no recursivas, y no a través de la recursión característica de los lenguajes funcionales clásicos. La repetición de un patrón sobre una estructura de tamaño conocido (por ejemplo, aplicar una puerta a cada cúbit de un registro) se expresa mediante combinadores de orden superior (como \lstinline|map|) cuyo número de pasos queda fijado al compilar, de modo que el circuito resultante es siempre finito y materializable. Esta decisión no es una carencia del lenguaje, sino el reflejo directo, en su modelo de cómputo, de lo que el sustrato físico puede ejecutar.

\section{El cúbit como recurso lineal}
\label{sec:modelo-lineal}

El núcleo semántico de Qalc consiste en tratar las variables cuánticas no como valores reutilizables, sino como recursos lineales, es decir, entidades que representan un circuito en construcción y que deben consumirse exactamente una vez.

\subsection{El problema del aliasing y la no clonación}

El teorema de no clonación establece que no existe ninguna operación física capaz de duplicar un estado cuántico arbitrario desconocido. En un lenguaje con semántica de valores convencional, la simple reutilización de una variable introduce aliasing (dos referencias al mismo estado). Una expresión tan inocua como usar \lstinline|q| dos veces sugeriría duplicar la información cuántica que \lstinline|q| representa, lo cual no tiene contrapartida física válida. Cualquier lenguaje cuántico de alto nivel debe impedir esta situación; la diferencia entre propuestas reside en cuándo y cómo la impide. Qalc opta por impedirla estáticamente.

\subsection{Variables cuánticas como circuitos en construcción}

La reconciliación con la inmutabilidad funcional se obtiene mediante un cambio de interpretación. Una variable de tipo cuántico no denota un estado físico mutable, sino el cable o la línea de un circuito en construcción. Aplicar una puerta no muta el estado: produce una nueva descripción de circuito que extiende la anterior. Como el valor original deja de estar disponible una vez transformado, no hay mutación observable y se preserva la transparencia referencial propia del paradigma funcional. El precio de esta interpretación es que el valor original debe quedar inaccesible tras su uso: de ahí la disciplina lineal.

\subsection{Disciplina lineal y consumo único}
\label{sec:disciplina_lineal}

Una variable cuántica debe usarse exactamente una vez a lo largo de su ámbito. No puede descartarse (es decir, usarse cero veces), porque ello equivaldría a perder un cúbit sin medir ni liberar, lo que constituye una fuga de recurso. Tampoco puede duplicarse (es decir, usarse dos o más veces), porque ello equivaldría a clonar información cuántica.

La linealidad es una propiedad intrínseca de los tipos cuánticos provistos por la biblioteca del lenguaje. Escribir código que respete la disciplina de consumo único de forma manual resulta verboso y repetitivo, pues obliga a re-enlazar explícitamente cada recurso tras cada operación. Para evitarlo, Qalc ofrece el bloque \lstinline|linear { ... }|, una construcción ergonómica que automatiza ese re-enlazado; se describe en detalle en la \autoref{sec:linear}.

\section{Sistema de tipos}

Qalc es un lenguaje fuertemente tipado. Su sistema de tipos combina tipos clásicos convencionales, tipos cuánticos abstractos de alto nivel, tipos función y una noción de linealidad que atraviesa todo el sistema. Conviene anticipar que la gramática del lenguaje no privilegia sintácticamente ningún tipo cuántico concreto, pues nombres como \lstinline|qbool| o \lstinline|qi16| son, a efectos del analizador sintáctico, identificadores de tipo ordinarios definidos por la biblioteca del lenguaje. Lo que distingue a un tipo cuántico de uno clásico es su carácter lineal.

\subsection{Tipos clásicos}

Los tipos clásicos no están sujetos a la disciplina lineal, por lo que pueden descartarse y duplicarse libremente. El lenguaje contempla los siguientes:

\begin{itemize}
    \item Enteros con anchura y signo explícitos. Los literales admiten los sufijos \lstinline|u8|, \lstinline|u16|, \lstinline|u32|, \lstinline|u64|, \lstinline|i8|, \lstinline|i16|, \lstinline|i32| e \lstinline|i64|, y pueden escribirse en base decimal, hexadecimal (\lstinline|0x...|) o binaria (\lstinline|0b...|). Un literal sin sufijo es un entero de anchura aún por determinar.
    \item Números en coma flotante con sufijos \lstinline|f32| y \lstinline|f64|, en notación decimal o exponencial.
    \item Tuplas, incluida la tupla vacía \lstinline|()|, escritas \lstinline|(a, b, c)|.
    \item Slices, tanto enumerados \lstinline|[1, 2, 3]| como por rango \lstinline|[1..10]|, y su tipo correspondiente \lstinline|[T]|.
\end{itemize}

El lenguaje no incorpora literales booleanos ni de cadena como categorías léxicas propias; los valores de verdad clásicos se modelan como constantes de un tipo definido en la biblioteca.

\subsection{Tipos cuánticos abstractos}

Sobre la base anterior, la biblioteca del lenguaje define tipos cuánticos abstractos de alto nivel. Entre ellos están \lstinline|qbool| (un cúbit lógico interpretado como valor booleano cuántico) y los enteros cuánticos con anchura y signo explícitos, como \lstinline|qi8| o \lstinline|qi16| (registros de 8 o 16 cúbits interpretados como enteros con signo), análogos cuánticos de los enteros clásicos. Todos estos tipos comparten una misma representación interna, pues son, internamente, un \lstinline|qcirc| (el circuito de cúbits), y se diferencian únicamente por la interpretación de alto nivel que cada uno impone sobre ese circuito. Todos son lineales por definición, de modo que toda variable que los habite queda sometida a la disciplina de consumo único.

El lenguaje proporciona además el tipo genérico \lstinline|qio<T>|, un tipo de entrada/salida cuántica que encapsula un cómputo con efectos sobre la QPU cuyo resultado es de tipo \lstinline|T|. Los valores \lstinline|qio| se componen dentro de bloques \lstinline|perform|.

\subsection{Tipos función y genéricos}

Los tipos función se escriben con flecha, \lstinline|A -> B|, asociando a la derecha, lo que habilita la aplicación parcial y el estilo currificado. Toda función de Qalc lleva asociada una anotación de tipo obligatoria de la forma \lstinline|nombre :: tipo|, al estilo de Haskell.

Los tipos genéricos (parametrizados) se escriben con parámetros entre ángulos, por ejemplo \lstinline|List<Int>|, lo que permite definir contenedores y abstracciones polimórficas.

\subsection{La disciplina lineal en el uso de los recursos cuánticos}
\label{sec:disciplina_lineal_uso}

El sistema de tipos divide los valores de un programa en dos categorías según cómo pueden usarse. Los valores clásicos (enteros, flotantes, booleanos clásicos, etc.) son ordinarios y pueden copiarse, reutilizarse cuantas veces se quiera o no usarse en absoluto. Los valores cuánticos (\lstinline|qbool|, \lstinline|qi16| y demás tipos cuánticos) son lineales, por lo que cada uno debe usarse exactamente una vez.

Conviene precisar qué significa usar un valor cuántico. Aplicarle una operación lo consume, de modo que a partir de ese momento ese valor deja de estar disponible. A cambio, la operación devuelve un valor cuántico nuevo (el resultado de la transformación), que el programa debe a su vez usar. Así, un cúbit fluye a través del programa como un testigo que se pasa de mano en mano: en cada paso hay una única copia viva, que se entrega a la siguiente operación.

De esta regla se siguen dos prohibiciones, cada una de las cuales evita un error físicamente imposible.

\paragraph{No se puede usar un valor cuántico dos veces (no clonación)}
Reutilizar una variable cuántica ya consumida equivaldría a duplicar información cuántica, algo que el teorema de no clonación prohíbe. El compilador rechaza el siguiente código:

\begin{lstlisting}[language=Qalc]
// ERROR: x se usa dos veces; seria clonar el cúbit.
duplicar :: qbool -> (qbool, qbool)
duplicar x = (H x, H x);
\end{lstlisting}

\paragraph{No se puede dejar un valor cuántico sin usar (no fuga)}
Recibir un cúbit y descartarlo dejaría un recurso físico sin medir ni liberar. El compilador también rechaza esto:

\begin{lstlisting}[language=Qalc]
// ERROR: x nunca se usa; el cúbit quedaria "colgado".
descartar :: qbool -> i16
descartar x = 0i16;
\end{lstlisting}

El uso correcto consume el recurso exactamente una vez y propaga el resultado:

\begin{lstlisting}[language=Qalc]
// Correcto: x se consume una vez y se devuelve el cúbit transformado.
aplicarH :: qbool -> qbool
aplicarH x = H x;
\end{lstlisting}

En contraste, un valor clásico no está sujeto a ninguna de estas restricciones. El siguiente código es válido, porque \lstinline|n| es un entero clásico y puede usarse varias veces o ninguna:

\begin{lstlisting}[language=Qalc]
// Correcto: n es clasico; se puede duplicar y descartar libremente.
doble :: i16 -> i16
doble n = n + n;
\end{lstlisting}

\paragraph{La linealidad se propaga a los tipos compuestos}
Una estructura que contenga, aunque sea indirectamente, algún valor cuántico hereda su carácter lineal, pues consumir o duplicar la estructura implicaría hacer lo propio con el recurso cuántico que encierra. Concretamente, una tupla es lineal en cuanto al menos uno de sus componentes lo es; así, \lstinline|(qi16, i16)| y \lstinline|(qbool, qbool)| son tipos lineales, mientras que \lstinline|(i16, i16)|, sin ningún componente cuántico, es clásico.

En la práctica, el compilador lleva la cuenta, para cada variable cuántica dentro de su ámbito, de si ya ha sido consumida. Si una variable se usa una segunda vez, o si su ámbito termina sin que se haya usado, se produce un error de compilación. Esta contabilidad es lo que convierte la no clonación y la ausencia de fugas en garantías estáticas; los programas que las violarían sencillamente no compilan, y el error se detecta antes de llegar a la QPU, no durante una ejecución costosa.

Aplicar manualmente esta disciplina obliga a re-enlazar el recurso tras cada operación, lo que resulta tedioso. El bloque \lstinline|linear| (véase la \autoref{sec:linear}) automatiza ese re-enlazado y es, por tanto, la forma idiomática de escribir secuencias de operaciones sobre recursos cuánticos sin incurrir en errores de uso.

\section{Léxico y estructura de un programa}

\subsection{Estructura general}

Un programa Qalc es una secuencia de declaraciones de nivel superior, esencialmente definiciones de funciones. No existe un punto de entrada implícito a nivel sintáctico; la organización del programa se basa enteramente en estas declaraciones.

\subsection{Espacios en blanco y comentarios}

Los espacios, tabuladores y saltos de línea son insignificantes y sirven únicamente como separadores. El lenguaje admite dos formas de comentario. Los comentarios de línea, introducidos por \lstinline|//| o por \lstinline|#|, se extienden hasta el final de la línea. Los comentarios de bloque, delimitados por \lstinline|/*| y \lstinline|*/|, pueden anidarse.

\subsection{Identificadores y palabras clave}

Un identificador simple comienza por una letra o por el carácter de subrayado y continúa con letras, dígitos o subrayados. Las siguientes palabras están reservadas y no pueden usarse como identificadores:

\lstinline|if|, \lstinline|else|, \lstinline|let|, \lstinline|type|, \lstinline|in|, \lstinline|then|, \lstinline|match|, \lstinline|perform|, \lstinline|linear|, \lstinline|where|, \lstinline|return|.

\subsection{Operadores}

La siguiente tabla resume los operadores del lenguaje agrupados por categoría.

\begin{table}[h]
\centering
\begin{tabularx}{\textwidth}{|l|X|}
\hline
\textbf{Categoría} & \textbf{Operadores} \\
\hline
Aplicación & \lstinline|$| (apply), \lstinline{|>} (pipe) \\
\hline
Lógicos & \lstinline{||}, \lstinline{&&} \\
\hline
Comparación & \lstinline|==|, \lstinline|!=|, \lstinline|<|, \lstinline|>|, \lstinline|<=|, \lstinline|>=| \\
\hline
Aritméticos & \lstinline|+|, \lstinline|-|, \lstinline|*|, \lstinline|/|, \lstinline|%|, \lstinline|**| (potencia) \\
\hline
A nivel de bit & \lstinline{|}, \lstinline|&|, \lstinline|^| \\
\hline
\end{tabularx}
\end{table}

\section{Expresiones y control de flujo}

El lenguaje carece de sentencias en el sentido imperativo, pues todo es una expresión que denota un valor. Se distinguen las expresiones atómicas, las expresiones de bloque y las expresiones binarias. Conviene recordar que ninguna de estas construcciones habilita la recursión, ya que la repetición se expresa mediante combinadores de orden superior acotados.

\subsection{Aplicación, lambdas y aplicación parcial}

La aplicación de funciones se realiza por yuxtaposición, \lstinline|f a b|, al estilo funcional. Las funciones anónimas (lambdas) se escriben \lstinline|(x y) -> x + y|, enumerando sus parámetros entre paréntesis. La aplicación parcial de operadores permite fijar uno de los operandos de un operador binario; así, \lstinline|(+ 5)| y \lstinline|(5 *)| denotan, respectivamente, la función que suma cinco y la que multiplica por cinco.

\subsection{Tuberías y aplicación explícita}

El operador de tubería \lstinline{|>} aplica una función en estilo fluido: \lstinline|a |> f| es equivalente a \lstinline|f a|, y encadena de forma legible una secuencia de transformaciones, lo que resulta especialmente natural para describir la composición de operaciones sobre un circuito.

\subsection{Bloques monádicos \texttt{perform} y \texttt{linear}}
\label{sec:linear}

Dos construcciones comparten la misma estructura de bloque monádico (una secuencia de enunciados separados por comas que puede terminar en una expresión \lstinline|return|). Por un lado, \lstinline|perform { ... }| secuencia operaciones con efectos, en el espíritu de la notación \textit{do} de Haskell. Por otro, \lstinline|linear { ... }| facilita el encadenado de operaciones sobre recursos lineales.

Dentro de estos bloques, además de expresiones y ligaduras ordinarias (\lstinline|p = e|), se admite la ligadura monádica \lstinline|p <- e|. Ambos bloques son azúcar sintáctico, pero se expanden de forma distinta, acorde con su propósito.

El bloque \lstinline|perform| secuencia cómputos con efectos (por ejemplo, valores \lstinline|qio|) y se expande a la combinación monádica habitual de las funciones \lstinline|flatMap| (encadenar un cómputo con efectos y ligar su resultado) y \lstinline|pure| (producir un cómputo sin efectos a partir de un valor).

El bloque \lstinline|linear|, por su parte, merece una explicación detallada, pues persigue un objetivo distinto, a saber, evitar la verbosidad del código que manipula variables lineales. Como una variable lineal se consume al usarse, encadenar varias operaciones sobre un mismo recurso de forma explícita obliga a re-enlazarlo tras cada paso, lo que resulta tedioso y propenso a errores. El bloque \lstinline|linear| automatiza ese re-enlazado, pues el compilador lo expande en una cadena anidada de expresiones \lstinline|let|, re-ligando automáticamente las variables lineales consumidas a sus mismos nombres.

\section{Operaciones cuánticas y composición de circuitos}

Qalc está pensado para escribirse, en la medida de lo posible, en un nivel de abstracción alto, donde el programador describe transformaciones sobre datos cuánticos (\lstinline|qbool|, \lstinline|qi16|, etc.) sin manipular cúbits físicos individuales. No obstante, el lenguaje también ofrece un nivel inferior para los casos en que es necesario controlar el circuito de forma explícita.

\subsection{Operaciones sobre datos cuánticos en el nivel alto}

En el nivel alto, las operaciones cuánticas no son construcciones sintácticas privilegiadas, sino funciones de la biblioteca del lenguaje con tipos lineales. Toman datos cuánticos, los consumen y devuelven los datos transformados. Por ejemplo, una operación que eleva un entero cuántico a una potencia constante, o que comprueba la igualdad con una constante produciendo un cúbit de resultado, tendrían signaturas de la forma siguiente.

\begin{lstlisting}[language=Qalc]
qPowConst :: i16 -> qi16 -> qi16
qEqConst  :: i16 -> qi16 -> (qi16, qbool)
\end{lstlisting}

La composición se expresa con las herramientas funcionales generales del lenguaje, donde la tubería \lstinline{|>} encadena transformaciones, las lambdas describen subcircuitos parametrizados y el bloque \lstinline|linear| secuencia operaciones que comparten recursos cuánticos respetando la disciplina lineal.

\subsection{El tipo \texttt{qcirc} en el nivel bajo}

Cuando se necesita programar el circuito explícitamente, se trabaja con el tipo \lstinline|qcirc|, que representa un circuito de $n$ cúbits. La convención de este nivel es que una puerta aplicada a un \lstinline|qcirc| se aplica a todos sus cúbits a la vez; al resto de los cúbits, cuando proceda, se les aplica implícitamente la identidad. Las puertas básicas son, por tanto, transformaciones del circuito completo.

\begin{lstlisting}[language=Qalc]
H :: qcirc -> qcirc
X :: qcirc -> qcirc
\end{lstlisting}

Para restringir una operación a cúbits concretos se emplea la función de orden superior \lstinline|at|, que recibe el índice sobre el que actuar y la puerta a aplicar. Está sobrecargada para aceptar tanto un índice único como una lista de índices.

\begin{lstlisting}[language=Qalc]
at :: i32 -> (qcirc -> qcirc) -> qcirc -> qcirc    // un cúbit
at :: [i32] -> (qcirc -> qcirc) -> qcirc -> qcirc  // varios cúbits

at 0 H                // aplica H solo al cúbit 0
at [0..3] X           // aplica X a los cúbits del 0 al 3
\end{lstlisting}

Para las puertas controladas se emplea la operación de orden superior \lstinline|controlled|.

\begin{lstlisting}[language=Qalc]
controlled :: i32 -> i32 -> (qcirc -> qcirc) -> qcirc -> qcirc

controlled 0 1 X      // CNOT con control en el cúbit 0 y objetivo el cúbit 1
\end{lstlisting}

Obsérvese que \lstinline|at 0 H|, \lstinline|at [0..3] X| y \lstinline|controlled 0 1 X| son aplicaciones parciales cuyo resultado es, en todos los casos, una función \lstinline|qcirc -> qcirc|. Esto las hace directamente componibles con la tubería \lstinline{|>}, lo que permite describir un circuito como una sucesión legible de transformaciones que se lee en el mismo orden en que se aplican las puertas. Por ejemplo, un circuito que prepara un estado de Bell sobre dos cúbits se escribe de la siguiente forma.

\begin{lstlisting}[language=Qalc]
// Estado de Bell sobre un circuito de 2 cúbits.
bell :: qcirc -> qcirc
bell c = c |> at 0 H |> controlled 0 1 X;
\end{lstlisting}

\subsection{Oráculo de Grover para $x^3 = 125$}

El algoritmo de búsqueda de Grover ilustra cómo el lenguaje combina tipos lineales, el bloque \lstinline|linear| y la composición funcional en el nivel alto. El elemento característico del algoritmo es el oráculo, es decir, una función que marca las soluciones del problema. Aquí se busca el conjunto de valores \lstinline|x| que satisfacen \lstinline|x³ = 125|. En Qalc, el oráculo es una función ordinaria.

\begin{lstlisting}[language=Qalc]
oracle :: qi16 -> (qi16, qbool)
oracle x = linear {
  cube <- qPowConst 3i16 x,         // x^3
  control <- qEqConst 125i16 cube,  // x^3 = 125
  uncompute cube,
  return (x, control)
}
\end{lstlisting}

En el oráculo se aprecian con claridad los rasgos del diseño. La variable cuántica \lstinline|x| entra una vez y se devuelve una vez, en la tupla final: nunca se duplica ni se descarta. El cómputo intermedio \lstinline|cube| es un recurso auxiliar: se calcula, se usa para producir el cúbit \lstinline|control| y, antes de devolver el resultado, se descomputa para liberar la memoria cuántica que ocupaba. El bloque \lstinline|linear| re-enlaza automáticamente los recursos entre pasos, de modo que la disciplina de uso único se respeta por construcción.

\chapter{Arquitectura del compilador}
\label{ch:arquitectura}

El presente capítulo describe la arquitectura interna del compilador de Qalc, en concreto la organización del código en crates, la abstracción compartida que conecta las fases del pipeline y el diseño de cada etapa de transformación. Para cada fase se describe qué representación recibe como entrada, qué transformación aplica y qué representación produce como salida, justificando las decisiones de diseño más relevantes.

\section{Visión general del pipeline}

El compilador de Qalc es un pipeline de transformaciones progresivas. Cada fase toma una representación del programa y produce otra más refinada, enriquecida con información que las fases siguientes necesitan. La cadena completa se muestra en la \autoref{tab:pipeline}.

\begin{table}[h]
\centering
\caption{Fases del pipeline del compilador de Qalc}
\label{tab:pipeline}
\small
\begin{tabularx}{\textwidth}{|l|l|X|}
\hline
\textbf{Fase} & \textbf{Crate / módulo} & \textbf{Transformación} \\
\hline
Análisis sintáctico & \lstinline|qalc-parser| & Texto fuente $\to$ \lstinline|ConcreteAST| \\
\hline
Normalización & \lstinline|qalc-analysis/normalizer| & \lstinline|ConcreteAST| $\to$ \lstinline|NormalizedAST| \\
\hline
Inferencia de tipos & \lstinline|qalc-analysis/hm_inferrer| & \lstinline|NormalizedAST| $\to$ \lstinline|TypedAST| \\
\hline
Verificación de recursividad & \lstinline|qalc-analysis| & \lstinline|TypedAST| $\to$ \lstinline|TypedAST| \\
\hline
Inliner vago & \lstinline|qalc-analysis| & \lstinline|TypedAST| $\to$ \lstinline|FlatAST| \\
\hline
Generación de código & \lstinline|qalc-codegen| & \lstinline|FlatAST| $\to$ OpenQASM~3 \\
\hline
\end{tabularx}
\end{table}

El punto de entrada ejecutable reside en el crate \lstinline|qalc-cli|, que encadena todas las fases y emite el resultado final por la salida estándar.

\section{El rasgo \texttt{CompilerStep} y la composición de fases}

Todas las fases del compilador comparten una interfaz común definida en \lstinline|qalc-utils|:

\begin{lstlisting}[language={}, caption=Definición del rasgo CompilerStep en Rust]
pub trait CompilerStep {
    type Input;
    type Output;
    fn transform(
        &self,
        input: Self::Input,
        diagnostics: &mut Vec<Diagnostic>,
    ) -> CompilerResult<Self::Output>;
}
\end{lstlisting}

El rasgo establece un contrato uniforme, en el que cada fase recibe su representación de entrada y un vector de diagnósticos acumulable, y devuelve \lstinline|CompilerResult<Output>|, que es \lstinline|Result<Output, Diagnostic>|. Los errores fatales cortan la ejecución del pipeline; los avisos se acumulan en el vector de diagnósticos y se reportan al final. Gracias a este contrato, añadir una nueva fase no requiere modificar las fases existentes; basta con definir un nuevo tipo que implemente \lstinline|CompilerStep| con la entrada y la salida apropiadas.

Cada diagnóstico lleva asociado un \lstinline|Span| (un par de posiciones en el texto fuente) que la biblioteca \texttt{ariadne}~\cite{ariadne} utiliza para renderizar mensajes de error con contexto de código fuente subrayado.

\section{Organización en crates}

El compilador se estructura como un workspace de Cargo con crates de dependencia estrictamente ordenada de abajo hacia arriba.

\begin{itemize}
  \item \texttt{qalc-utils.} Es la base de la pila. Define \lstinline|Source| (el texto fuente con su nombre de fichero), \lstinline|Span| (intervalo de bytes dentro de un \lstinline|Source|), \lstinline|Diagnostic| (un error o aviso con span y mensaje) y el rasgo \lstinline|CompilerStep|. Ningún otro módulo puede ignorar un span ni construir un error sin una ubicación en el código fuente.

  \item \texttt{qalc-ast.} Define todos los tipos de nodo AST usados a lo largo del pipeline: \lstinline|ConcreteAST|, \lstinline|NormalizedAST|, \lstinline|TypedAST| y los módulos auxiliares de ámbito léxico (\lstinline|scope|) y metadatos de tipo (\lstinline|typed_ast|). Al centralizar los tipos de nodo se evita la proliferación de estructuras de datos redundantes entre crates.

  \item \texttt{qalc-builtins.} Expone el tipo \lstinline|BuiltinDef| y la función \lstinline|standard_builtins()|, que devuelve las definiciones de todas las funciones predefinidas del lenguaje (operaciones aritméticas, comparaciones, etc.) junto con su tipo. Este crate desacopla la definición de los built-ins de la implementación de la inferencia de tipos.

  \item \texttt{qalc-parser.} Implementa el analizador léxico-sintáctico. Depende únicamente de \lstinline|qalc-utils| y \lstinline|qalc-ast|.

  \item \texttt{qalc-analysis.} Contiene el normalizador y el inferidor de tipos HM. Depende de los cuatro crates anteriores.

  \item \texttt{qalc-cli.} Es el ejecutable. Importa todas las fases, las encadena y presenta los resultados.
\end{itemize}

\section{El patrón AST paramétrico}
\label{sec:ast-parametrico}

Una decisión de diseño transversal del compilador es que los tipos de nodo del AST son genéricos en sus metadatos. Cada nodo de expresión, declaración, patrón o tipo lleva un parámetro de tipo \lstinline|M: ASTMetadata| que determina qué información está adjunta a cada nodo en esa fase del pipeline:

\begin{lstlisting}[language={}, caption=Definición del rasgo ASTMetadata en Rust]
pub trait ASTMetadata {
    type DeclarationMeta;
    type ExpressionMeta;
    type PatternMeta;
    type TypeMeta;
}
\end{lstlisting}

El crate \lstinline|qalc-ast| define dos implementaciones de este rasgo que se emplean a lo largo del pipeline. La implementación \lstinline|Spanned| (en la que todos los campos meta son \lstinline|Span|) se usa en el AST concreto y en el AST normalizado. La implementación \lstinline|Typed| (en la que los campos meta son \lstinline|TypeMetadata \{ span, ttype, scope \}| con el tipo completamente resuelto) se usa en el AST tipado.

Existe además \lstinline|Inferred| (con \lstinline|TypeWithUnknow| en lugar de \lstinline|Type|), que se usa internamente durante la inferencia de tipos antes de que la unificación resuelva todas las variables de tipo. La conversión de \lstinline|Inferred| a \lstinline|Typed| ocurre en la fase de aplicación de sustituciones descrita en el \autoref{ch:analisis-semantico}.

Gracias a este patrón, los tipos de nodo se definen una sola vez y se reutilizan sin modificación en todas las fases. Añadir una nueva fase, por ejemplo una que propague información de linealidad, implica únicamente definir un nuevo \lstinline|ASTMetadata| sin tocar los nodos existentes.

\section{Análisis sintáctico}
\label{sec:parser}

El analizador sintáctico transforma el texto fuente en un \lstinline|ConcreteAST|. Se articula en dos capas.

\subsection{Gramática PEG con Pest}

La gramática del lenguaje se define en un único fichero \lstinline|grammar.pest| usando la biblioteca Pest~\cite{pest_parser}, que implementa gramáticas de expresión de análisis (PEG, del inglés \textit{Parsing Expression Grammar})~\cite{ford_peg_2004}. Las PEG son deterministas por construcción: no hay ambigüedad, cada alternativa se intenta en orden y la primera que encaja gana. Esto elimina la ambigüedad habitual de las gramáticas libres de contexto y simplifica el mantenimiento.

El fichero \lstinline|grammar.pest| es la fuente de verdad para la sintaxis del lenguaje. Cualquier nueva construcción sintáctica debe reflejarse primero en la gramática y después en el analizador; la dirección inversa no es válida.

\subsection{Análisis de Pratt para operadores}

Las expresiones aritméticas, lógicas y de comparación se analizan con un analizador de Pratt~\cite{pratt_1973}, también llamado \textit{top-down operator precedence}. Este algoritmo asocia a cada operador un poder de enlace (\textit{binding power}) izquierdo y derecho, y construye el árbol de expresiones respetando la precedencia y la asociatividad sin necesidad de reglas gramaticales adicionales. El resultado es una representación de expresión binaria encadenada que el normalizador transforma en aplicaciones currificadas.

\section{Normalización}
\label{sec:normalizador}

El normalizador (\lstinline|qalc-analysis/normalizer|) convierte el \lstinline|ConcreteAST| en un \lstinline|NormalizedAST| más simple. La representación normalizada elimina construcciones sintácticas de alto nivel dejando un núcleo reducido que las fases siguientes pueden procesar de manera uniforme.

\begin{itemize}
  \item \textit{Desazucarado de llamadas a función.} Las llamadas a función con múltiples argumentos, de la forma \lstinline|f a b c|, se transforman en una cadena de aplicaciones binarias currificadas.

\begin{lstlisting}[language={}, caption=Desazucarado de llamadas a función con múltiples argumentos]
-- Fuente:          f a b c
-- Normalizado:     ((f a) b) c
\end{lstlisting}

  En el NormalizedAST solo existe un nodo \lstinline|Transform \{ callee, param \}| que representa la aplicación de una función a un único argumento. Esto simplifica el inferidor de tipos, que no necesita distinguir entre funciones variadicas y funciones unarias.

  \item \textit{Desazucarado de patrones en parámetros.} Qalc permite deconstruir tuplas directamente en los parámetros de una función o lambda. El normalizador elimina ese azúcar sintáctico reemplazando el patrón por un nombre generado y añadiendo un bloque \lstinline|let| al inicio del cuerpo.

\begin{lstlisting}[language={}, caption=Desazucarado de deconstrucción de tupla en parámetros]
-- Fuente:
add (x, y) = x + y;

-- Normalizado:
add p0 = let (x, y) = p0 in x + y;
\end{lstlisting}

  \item \textit{Desazucarado de} \texttt{match}. Las expresiones \lstinline|match| se transforman en cadenas de \lstinline|if-else| anidados. Cada rama del \lstinline|match| se convierte en una comparación de igualdad con el valor escrutado; el caso por defecto pasa a ser la alternativa del último \lstinline|else|. Esta simplificación evita introducir un nodo \lstinline|Match| en el inferidor y en las fases siguientes.

\begin{lstlisting}[language={}, caption=Desazucarado de match a cadena de if-else]
-- Fuente:
match x { 1 -> "uno", 2 -> "dos" }

-- Normalizado:
if x == 1 then "uno" else if x == 2 then "dos" else ()
\end{lstlisting}

  \item \textit{Desazucarado de bloques} \texttt{linear} \textit{y} \texttt{perform}. Ambos bloques monádicos se desazucaran en el normalizador. El bloque \lstinline|linear|, que facilita el encadenamiento de operaciones sobre recursos lineales, se expande en una cadena de expresiones \lstinline|let| anidadas que re-enlazan automáticamente las variables lineales consumidas. La ligadura monádica \lstinline|x <- f a| extrae el resultado del par devuelto por \lstinline|f a| y lo asigna a \lstinline|x|, relanzando la continuación con los valores re-enlazados. El bloque \lstinline|perform| sigue un esquema análogo orientado a la composición de efectos.
\end{itemize}

\section{Verificador de recursividad}
\label{sec:recursivity-checker}

La fase de verificación de recursividad recorre el \lstinline|TypedAST| comprobando que ninguna declaración de función se invoca a sí misma directa o indirectamente. Esta restricción no es arbitraria, pues un circuito cuántico es una estructura finita y estática que debe quedar completamente determinada en tiempo de compilación, y OpenQASM~3 no dispone de llamadas recursivas ni de una pila dinámica de invocaciones~\cite{openqasm3_2022}. Permitir recursividad generaría un circuito de profundidad potencialmente infinita, lo que carece de contrapartida física ejecutable.

El diseño de la fase consiste en construir el grafo de llamadas del programa a partir del \lstinline|TypedAST|, donde cada declaración de función es un nodo y existe una arista dirigida de $f$ a $g$ si el cuerpo de $f$ contiene alguna referencia al nombre $g$. Un ciclo en este grafo indica recursividad. El algoritmo empleado para detectar ciclos es una búsqueda en profundidad (DFS, del inglés \textit{Depth-First Search}) con marcado de nodos en tres estados (no visitado, en pila, completado), una variante clásica del algoritmo de Tarjan para componentes fuertemente conexas~\cite{aho_compiladores_2006}. Cuando se detecta un nodo en pila al explorar sus vecinos, se tiene un ciclo; el verificador recupera el span del identificador que cierra el ciclo y emite un error de compilación con el camino completo de llamadas para ayudar al programador a localizar la recursividad.

Esta fase actúa como una validación pura, es decir, no modifica el AST ni produce una representación nueva. Su única salida observable es la aprobación o el rechazo del programa con el diagnóstico correspondiente. Al situarse después de la inferencia de tipos, el grafo de llamadas ya tiene información de tipos resuelta, lo que permitirá en el futuro extender la verificación para detectar también ciclos de valor mutuamente dependientes en el sentido de la evaluación vaga.

\section{Inliner vago}
\label{sec:inliner}

El inliner vago transforma el \lstinline|TypedAST| en un \lstinline|FlatAST|, una representación aplanada orientada directamente a la generación de código. La motivación es que el código OpenQASM~3 que se emitirá en la fase siguiente no tiene noción de abstracción funcional, pues es una secuencia plana de instrucciones de puerta sobre registros concretos. El inliner cierra esa brecha materializando, en tiempo de compilación, todas las aplicaciones de función.

El mecanismo es el de una evaluación vaga (\textit{lazy evaluation}) simbólica: el inliner recorre el AST como si fuera un intérprete, pero en lugar de producir valores concretos, produce nodos del \lstinline|FlatAST|. Las expresiones cuyo valor puede determinarse completamente en tiempo de compilación (literales, aplicaciones de funciones conocidas sobre argumentos conocidos) se reducen a su forma normal. Las expresiones que dependen de parámetros cuánticos de entrada, cuyo valor solo existe en tiempo de ejecución cuántica, se envuelven en thunks, es decir, nodos suspendidos del \lstinline|FlatAST| que representan un cómputo diferido y que el codegen traducirá directamente a secuencias de puertas.

La estrategia vaga es necesaria porque el lenguaje usa combinadores de orden superior como \lstinline|map| para expresar la aplicación repetida de una operación a un slice de cúbits. El inliner despliega estos combinadores en tiempo de compilación conociendo la longitud del slice (que es estática, ya que el tamaño de los registros cuánticos es fijo), produciendo una secuencia plana de aplicaciones de puerta sin ninguna abstracción pendiente.

El \lstinline|FlatAST| resultante contiene únicamente: referencias a built-ins con sus argumentos concretos, composiciones secuenciales de operaciones y condicionales clásicos basados en resultados de medición. No existe ningún nodo lambda, ningún bloque \lstinline|let| ni ninguna referencia a funciones definidas por el usuario: todas han sido inlinadas y sus cuerpos sustituidos en el punto de llamada.

\section{Generación de código}
\label{sec:codegen}

La fase de generación de código toma el \lstinline|FlatAST| producido por el inliner y emite texto OpenQASM~3 válido. OpenQASM~3 es el estándar abierto de descripción de circuitos cuánticos, adoptado por IBM Quantum, IonQ y Rigetti, entre otros proveedores~\cite{openqasm3_2022}. A diferencia de versiones anteriores, OpenQASM~3 incorpora construcciones de control de flujo clásico (condicionales y bucles acotados basados en resultados de medición) que se corresponden directamente con los nodos del \lstinline|FlatAST|.

El diseño del codegen se estructura en dos partes. El preámbulo del fichero de salida declara la versión de OpenQASM~3, los registros cuánticos y clásicos necesarios, y las definiciones de las puertas derivadas de los built-ins que no están incluidas en el conjunto estándar de OpenQASM~3. Cada \lstinline|BuiltinDef| del crate \lstinline|qalc-builtins| tiene asociada una plantilla de emisión: para las operaciones aritméticas clásicas se emiten expresiones de OpenQASM~3; para las puertas cuánticas básicas (Hadamard, Pauli-X, CNOT, etc.) se emiten las instrucciones de puerta estándar del lenguaje.

El cuerpo del circuito se genera recorriendo el \lstinline|FlatAST| en orden topológico y emitiendo para cada nodo la instrucción OpenQASM~3 correspondiente: aplicaciones de puerta sobre cúbits concretos, mediciones con asignación a bits clásicos, y bloques \lstinline|if| condicionados a resultados de medición para el control de flujo clásico-cuántico. La correspondencia entre nodos del \lstinline|FlatAST| e instrucciones OpenQASM~3 es, por diseño, casi uno a uno; el trabajo de aplanar y especializar el programa ya ha sido realizado por el inliner, de modo que el codegen puede limitarse a traducir sin necesidad de realizar transformaciones adicionales.

\chapter{Análisis semántico y sistema de tipos}
\label{ch:analisis-semantico}

El análisis semántico de Qalc tiene como objetivo dotar a cada nodo del AST de un tipo completamente resuelto, verificando que el programa es coherente desde el punto de vista tipológico antes de que ninguna fase posterior lo transforme en código ejecutable. Este proceso se apoya en dos mecanismos que se ejecutan de forma coordinada, el ámbito léxico, que resuelve a qué declaración se refiere cada identificador, y el inferidor de tipos basado en el algoritmo de Hindley-Milner, que asigna un tipo a cada expresión a partir de esas resoluciones. Este capítulo describe primero la resolución del ámbito léxico, después el algoritmo de inferencia con sus tres fases, la representación interna de los tipos durante la inferencia, el mecanismo de unificación en detalle y la sobrecarga de funciones discriminada por tipo.

\section{Ámbito léxico y declaraciones de nivel superior}
\label{sec:ambito-lexico}

El primer paso del análisis semántico consiste en resolver el ámbito léxico de las declaraciones, de modo que la fase de inferencia pueda tratar cada identificador exactamente como el símbolo al que se refiere. El módulo \lstinline|qalc-ast/scope.rs| implementa el ámbito léxico del compilador. Un \lstinline|Scope| es un mapa de nombres a vectores de \lstinline|SymbolInfo|, con un puntero opcional al ámbito padre para implementar el encadenamiento léxico:

\begin{lstlisting}[language={}, caption=Definición del tipo Scope para el ámbito léxico]
pub struct Scope {
    parent: Option<Rc<Scope>>,
    symbols: RefCell<HashMap<String, Vec<SymbolInfo>>>,
}
\end{lstlisting}

Hay dos operaciones de inserción con semánticas distintas. La primera, \lstinline|define(&self, symbol)|, registra un nombre nuevo y, si el nombre ya existe en el ámbito local, devuelve \lstinline|false| sin sobreescribirlo (se usa para parámetros y variables \lstinline|let|, donde la redefinición en el mismo ámbito es un error). La segunda, \lstinline|define_decl(&self, symbol)|, añade la definición al vector del nombre sin rechazar las entradas previas, lo que permite acumular varias definiciones para el mismo nombre (se usa para declaraciones de función de nivel superior, donde el mismo nombre puede tener varias definiciones con firmas de tipo distintas; la resolución de la sobrecarga se describe en la \autoref{sec:overloading}).

Al inicio de la inferencia, el \lstinline|HMInferrer| realiza dos operaciones preparatorias sobre el ámbito global antes de visitar ninguna declaración del programa.

En primer lugar, inyecta todas las funciones predefinidas del crate \lstinline|qalc-builtins|. Cada \lstinline|BuiltinDef| contiene el nombre de la función y su tipo completo expresado como \lstinline|TypeWithUnknow|, de modo que el inferidor puede resolver referencias a built-ins exactamente igual que referencias a funciones definidas por el usuario.

En segundo lugar, recorre todas las declaraciones del programa y registra sus firmas de tipo en el ámbito global antes de inferir sus cuerpos. Este paso de predeclaración es fundamental para que las funciones puedan referenciarse mutuamente sin que el orden de aparición en el fichero fuente afecte a la corrección de la inferencia.

\section{El algoritmo de Hindley-Milner}
\label{sec:hm}

Una vez resuelto el ámbito léxico de las declaraciones, Qalc infiere el tipo de cada expresión mediante el algoritmo de Hindley-Milner, que resuelve el problema de la inferencia de tipos en un lenguaje funcional sin necesidad de que el programador anote los tipos de cada expresión. A partir de las anotaciones de las declaraciones de nivel superior y de las reglas del sistema de tipos, el compilador deduce el tipo de toda subexpresión.

El nombre del algoritmo combina las aportaciones de dos autores. Hindley estudió en 1969 la asignación de tipos principales en el marco de la lógica combinatoria~\cite{hindley_1969}, y Milner reformuló esa idea en 1978 como un algoritmo de inferencia práctico, conocido como algoritmo W~\cite{milner_1978}, pensado para el lenguaje ML. Damas y Milner~\cite{damas_milner_1982} extendieron después el algoritmo para admitir polimorfismo en las declaraciones \lstinline|let|, dando lugar a la versión que se usa habitualmente hoy en día y que sirve de referencia para la implementación de Qalc.

El algoritmo de Hindley-Milner es la base de la inferencia de tipos en toda la familia de lenguajes ML (del inglés \textit{Meta Language}), que incluye Standard ML, OCaml y F#, y en Haskell, que lo extiende con clases de tipos para resolver la sobrecarga. Lenguajes más recientes como Elm o PureScript también lo adoptan sin modificaciones sustanciales. Otros lenguajes con un sistema de tipos más flexible, como Rust~\cite{matsakis_rust_2014}, Swift, Kotlin o TypeScript, usan variantes locales o bidireccionales inspiradas en Hindley-Milner que infieren el tipo de expresiones concretas sin llegar al polimorfismo global de la versión completa del algoritmo. Qalc adopta Hindley-Milner porque permite prescindir de anotaciones de tipo en la mayoría de las expresiones sin renunciar a la seguridad de tipos, una propiedad especialmente valiosa en un lenguaje pensado para describir circuitos cuánticos de forma concisa.

El algoritmo se estructura en tres fases que se ejecutan secuencialmente~\cite{pierce_tapl_2002}:

\begin{enumerate}
  \item Generación de restricciones. Se recorre el AST y a cada subexpresión se le asigna una variable de tipo fresca. A medida que se visitan los nodos, se generan restricciones de tipo, es decir, aserciones de la forma $\tau_1 = \tau_2$ que afirman que dos tipos deben ser iguales. Por ejemplo, al analizar la aplicación $f\;a$, se genera la restricción $\tau_f = \tau_a \to \tau_{resultado}$.
  \item Unificación. Las restricciones generadas se resuelven mediante el algoritmo de unificación de Robinson~\cite{robinson_1965}. El resultado es una sustitución, es decir, una función finita que asigna a cada variable de tipo fresca el tipo concreto (o el tipo estructurado) que la resuelve.
  \item Aplicación de la sustitución. La sustitución se aplica a todo el AST, reemplazando cada variable de tipo por su valor en la sustitución. El resultado es el \lstinline|TypedAST|, en el que cada nodo lleva su tipo completamente resuelto.
\end{enumerate}

En Qalc, estas tres fases se implementan en los módulos \lstinline|infer.rs|, \lstinline|unify.rs| y \lstinline|apply.rs| del crate \lstinline|qalc-analysis/hm_inferrer|, respectivamente.

Para ilustrar cómo cooperan las tres fases sobre un caso concreto, considérese la siguiente declaración.

\begin{lstlisting}[language={}, caption=Ejemplo de declaración a inferir mediante Hindley-Milner]
incrementa :: i32 -> i32
incrementa = \x -> x + 1
\end{lstlisting}

En la fase de generación de restricciones (\autoref{sec:constraint-gen}), la lambda asigna una variable fresca $\tau_x$ al parámetro \lstinline|x| y visita el cuerpo \lstinline|x + 1|. Como \lstinline|+| es una función predefinida de tipo \lstinline|i32 -> i32 -> i32|, la aplicación genera las restricciones $\tau_x = \mathtt{i32}$ y $\tau_{cuerpo} = \mathtt{i32}$, y la anotación de la declaración añade además la restricción $\tau_x \to \tau_{cuerpo} = \mathtt{i32} \to \mathtt{i32}$.

En la fase de unificación (\autoref{sec:unificacion}), estas restricciones se resuelven en una única sustitución que asigna $\tau_x \mapsto \mathtt{i32}$ y $\tau_{cuerpo} \mapsto \mathtt{i32}$, sin que se detecte ningún conflicto entre las variables y los tipos concretos involucrados.

En la fase de aplicación de la sustitución, esa sustitución se propaga a cada nodo de la lambda, de modo que el \lstinline|TypedAST| resultante registra que \lstinline|x| tiene tipo \lstinline|i32| y que \lstinline|incrementa| tiene tipo \lstinline|i32 -> i32|, coherente con la anotación original de la declaración.

\section{Representación de tipos durante la inferencia}

El crate \lstinline|qalc-ast| define dos enumeraciones para representar los tipos en distintos momentos del pipeline.

\begin{itemize}
  \item \texttt{TypeWithUnknow.} Es la representación usada durante la inferencia. Sus variantes son las siguientes: \lstinline|Unknow(usize)| (una variable de tipo fresca, identificada por un entero único generado con \lstinline|gen_id()|, que representa tipos aún no resueltos), \lstinline|Builtin \{ id, params, linear \}| (un tipo conocido, como \lstinline|i32|, \lstinline|bool| o \lstinline|qbool|, donde el campo \lstinline|linear| distingue los tipos cuánticos de los clásicos), \lstinline|Arrow(a, b)| (el tipo función $a \to b$), \lstinline|Tuple(ts)| (el tipo tupla $(t_1, t_2, \ldots, t_n)$) y \lstinline|Slice(t)| (el tipo slice $[t]$).

  \item \texttt{Type.} Es la representación final, producida una vez que la unificación ha resuelto todas las variables. Es isomorfa a \lstinline|TypeWithUnknow| salvo por la ausencia de la variante \lstinline|Unknow|: si la unificación no pudo resolver alguna variable, la conversión \lstinline|TryFrom<TypeWithUnknow> for Type| falla con un error de compilación que pide al usuario una anotación de tipo explícita.
\end{itemize}

La separación entre \lstinline|TypeWithUnknow| y \lstinline|Type| garantiza, a nivel del sistema de tipos de Rust, que ningún nodo del \lstinline|TypedAST| final pueda contener una variable sin resolver.

\section{Generación de restricciones de tipo}
\label{sec:constraint-gen}

La función \lstinline|infer_expr| recorre el AST normalizado en profundidad y produce un par \lstinline|(Expression<Inferred>, Vec<TypeConstraint>)|: el nodo anotado con su variable de tipo y el conjunto de restricciones que impone la estructura de ese nodo.

Las reglas más relevantes son las siguientes:

\begin{itemize}
  \item \textit{Literales.} Los literales enteros y flotantes producen directamente el tipo concreto determinado por su sufijo: \lstinline|42i32| produce \lstinline|Builtin("i32")|, \lstinline|3.14f64| produce \lstinline|Builtin("f64")|. Los literales sin sufijo explícito se tipan por defecto como \lstinline|i32| y \lstinline|f64| respectivamente.

  \item \textit{Identificadores.} La resolución de un identificador consulta el ámbito léxico en vigor. Si el nombre se encuentra, el nodo hereda el tipo registrado en el símbolo. Si no existe, se emite un error de compilación con el span del identificador.

  \item \textit{Aplicación de función} (\texttt{Transform \{ callee, param \}}). Se genera una variable fresca $\tau_{ret}$ para el tipo de retorno y se añade la restricción $\tau_{callee} = \tau_{param} \to \tau_{ret}$. El tipo del nodo de aplicación es $\tau_{ret}$.

  \item \textit{Lambda.} Se crea un ámbito nuevo, se registra cada parámetro con una variable fresca y se infiere el cuerpo. El tipo de la lambda se construye encadenando los tipos de los parámetros con el tipo del cuerpo, tal y como muestra la expresión
  \[\tau_{lambda} = \tau_{p_1} \to \tau_{p_2} \to \cdots \to \tau_{cuerpo}\]

  \item \textit{Bloque} \texttt{let}. Cada enlace $p = e$ genera las restricciones de $e$ y añade la restricción $\tau_e = \tau_p$, ligando el tipo de la expresión al tipo del patrón. Después se infiere la expresión de valor en el ámbito extendido. El tipo del bloque es el tipo de la expresión de valor.

  \item \textit{If-else.} Se genera una variable fresca $\tau_{ret}$ y se añaden las restricciones $\tau_{ret} = \tau_{then}$ y $\tau_{ret} = \tau_{else}$, exigiendo que ambas ramas tengan el mismo tipo. El tipo del test no se restringe a ningún tipo concreto en esta fase.

  \item \textit{Tuplas.} El tipo de una tupla $(e_1, \ldots, e_n)$ es \lstinline|Tuple|$(\tau_{e_1}, \ldots, \tau_{e_n})$, donde cada $\tau_{e_i}$ es el tipo inferido de la subexpresión correspondiente.

  \item \textit{Slices.} Se genera una variable fresca $\tau_{elem}$ y se añade la restricción $\tau_{elem} = \tau_{e_i}$ para cada elemento, exigiendo homogeneidad del slice.
\end{itemize}

La función \lstinline|infer_decl| maneja las declaraciones de función. Cada declaración tiene una anotación de tipo explícita escrita por el programador. El inferidor convierte esa anotación en un \lstinline|TypeWithUnknow| mediante \lstinline|type_expr_to_infer_type| y añade la restricción:
\[\tau_{inferida} = \tau_{anotada}\]

Esto obliga al cuerpo de la función a ser coherente con la firma declarada. Si la inferencia produce un tipo incompatible, el error se reporta apuntando al span de la declaración.

\section{Unificación}
\label{sec:unificacion}

La unificación resuelve el conjunto de restricciones generado en la fase anterior. El algoritmo implementado en \lstinline|unify.rs| es el de Robinson~\cite{robinson_1965}, en su variante con verificación de ocurrencias (\textit{occurs check}).

El algoritmo mantiene una sustitución parcial $S$ (un mapa de identificadores de variables a tipos). Para cada restricción $\tau_1 = \tau_2$, primero aplica $S$ a ambos lados (reemplazando las variables ya resueltas) y después llama a \lstinline|unify_one|:

\begin{itemize}
  \item Si $\tau_1 = \tau_2$, no hay nada que hacer.
  \item Si alguno de los dos lados es una variable fresca $?n$, se extiende $S$ con $n \mapsto \tau_{otro}$, siempre que $n$ no aparezca en $\tau_{otro}$ (occurs check). La verificación de ocurrencias previene la generación de tipos circulares infinitos como $\tau = \tau \to \tau$.
  \item Si ambos son tipos estructurados del mismo constructor (\lstinline|Arrow|, \lstinline|Tuple|, \lstinline|Slice|, \lstinline|Builtin|), se unifican recursivamente sus subcomponentes.
  \item En cualquier otro caso, los tipos son incompatibles y se emite un error de compilación con el span de la restricción original, indicando qué tipos no pudieron unificarse.
\end{itemize}

La implementación aplica la sustitución parcial acumulada antes de cada unificación, lo que permite que el resultado de resolver una restricción temprana informe de inmediato a las restricciones posteriores, sin necesidad de múltiples pasadas.

\section{Aplicación de la sustitución}

Una vez que \lstinline|unify| devuelve la sustitución final $S$, el módulo \lstinline|apply.rs| recorre el AST \lstinline|Inferred| de forma estructural, reemplazando cada \lstinline|TypeWithUnknow::Unknow(id)| por \lstinline|S(id)| en cada nodo. El resultado es el \lstinline|TypedAST|, en el que todas las variables de tipo han sido resueltas a tipos concretos y el AST puede convertirse al tipo final \lstinline|Typed|.

Si alguna variable permanece irresuelta tras la sustitución (por ejemplo, en una función sin anotación de tipo y con un polimorfismo no instanciado), la conversión \lstinline|TryFrom<TypeWithUnknow> for Type| falla y se reporta un error pidiendo al usuario que añada una anotación explícita.

\section{Sobrecarga de funciones discriminada por tipo}
\label{sec:overloading}

Una de las extensiones del sistema de tipos es la sobrecarga de funciones basada en el tipo del primer parámetro. En Qalc, una función puede tener varias definiciones con el mismo nombre siempre que difieran en el tipo del primer parámetro; el compilador elige la definición adecuada en tiempo de compilación consultando el tipo del primer argumento en el punto de llamada.

Este mecanismo permite expresar operaciones polimórficas sin necesidad de clases de tipos ni de polimorfismo paramétrico completo, manteniendo la inferencia de tipos decidible y eficiente. Por ejemplo, la función \lstinline|at|, que aplica una puerta a un cúbit concreto, puede estar definida dos veces con el mismo nombre:

\begin{lstlisting}[language={}, caption=Ejemplo de sobrecarga de función discriminada por tipo del primer parámetro]
at :: i32 -> (qcirc -> qcirc) -> qcirc -> qcirc
at :: [i32] -> (qcirc -> qcirc) -> qcirc -> qcirc
\end{lstlisting}

La sobrecarga está soportada a nivel de ámbito: \lstinline|define_decl| acumula todas las definiciones de un mismo nombre en un vector, y la operación \lstinline|resolve| recorre ese vector dado el tipo del primer argumento en el punto de llamada, y elige la definición cuyo primer tipo de parámetro es unificable con el tipo del argumento. Si ninguna definición es compatible, se emite un error de compilación indicando las firmas disponibles y el tipo que se intentó usar.

La discriminación ocurre tras la generación de restricciones para el primer argumento, cuando su tipo ya está disponible (posiblemente aún como variable fresca parcialmente restringida). En los casos en que el tipo del primer argumento no está suficientemente determinado en ese punto, el compilador difiere la elección hasta después de la unificación y resuelve la sobrecarga en una segunda pasada sobre el AST, manteniendo en el nodo de identificador una referencia al conjunto de candidatos hasta que la sustitución lo resuelve.

Esta estrategia de sobrecarga es deliberadamente más sencilla que los sistemas de clases de tipos de Haskell o los rasgos de Rust~\cite{matsakis_rust_2014}, ya que restringir la discriminación al tipo del primer parámetro garantiza que la resolución es siempre unívoca y que no es necesario un mecanismo de inferencia global adicional. El coste es una menor expresividad, que en el contexto de Qalc es aceptable dado que el objetivo del lenguaje es describir circuitos cuánticos, no implementar abstracciones de propósito general.

\section{Casos de uso}
\label{sec:casos-uso}

Los siguientes ejemplos recorren el proceso completo de análisis semántico descrito en este capítulo sobre fragmentos de código más cercanos al uso real del lenguaje que los ejemplos aislados de las secciones anteriores, mostrando cómo cooperan la resolución de ámbito, la inferencia de tipos, la unificación y la sobrecarga sobre un mismo programa.

\subsection{Inferencia sobre una operación cuántica compuesta}

Considérese la siguiente declaración, que aplica la puerta Hadamard a dos cúbits de entrada y devuelve ambos resultados en una tupla.

\begin{lstlisting}[language={}, caption=Declaración que combina un bloque let y una tupla sobre valores cuánticos]
prepara :: qbool -> qbool -> (qbool, qbool)
prepara a b = let x = H a in (x, H b)
\end{lstlisting}

La resolución de ámbito (\autoref{sec:ambito-lexico}) registra los parámetros \lstinline|a| y \lstinline|b| con sendas variables frescas $\tau_a$ y $\tau_b$, y resuelve el identificador \lstinline|H| contra la firma \lstinline|qbool -> qbool| que \lstinline|qalc-builtins| inyectó en el ámbito global.

Durante la generación de restricciones (\autoref{sec:constraint-gen}) se combinan cuatro de las reglas descritas en ese apartado. La aplicación \lstinline|H a| impone, según la regla de aplicación de función, la restricción $\mathtt{qbool} \to \mathtt{qbool} = \tau_a \to \tau_{ret_1}$, de la que se derivan $\tau_a = \mathtt{qbool}$ y $\tau_{ret_1} = \mathtt{qbool}$. La regla del bloque \lstinline|let| liga entonces $\tau_x = \tau_{ret_1}$. La segunda aplicación \lstinline|H b| genera de forma análoga $\tau_b = \mathtt{qbool}$ y una variable fresca $\tau_{ret_2} = \mathtt{qbool}$, y la regla de tuplas construye el tipo \lstinline|Tuple|$(\tau_x, \tau_{ret_2})$ para el resultado. Por último, la anotación de la declaración añade la restricción de que el tipo completo coincida con \lstinline|qbool -> qbool -> (qbool, qbool)|.

La unificación (\autoref{sec:unificacion}) resuelve todas las variables anteriores a \lstinline|qbool| sin ningún conflicto, y la sustitución confirma que \lstinline|prepara| tiene exactamente el tipo anotado. El carácter lineal de \lstinline|qbool| viaja con el tipo sin intervención adicional del inferidor, puesto que forma parte del propio campo \lstinline|linear| de la variante \lstinline|Builtin|; comprobar que \lstinline|a| y \lstinline|b| se consumen exactamente una vez es responsabilidad de la fase de comprobación de consumo lineal, posterior al aplanado y descrita en el \autoref{sec:linear-checker}.

\subsection{Resolución de sobrecarga en una llamada concreta}

Retomando la función \lstinline|at| introducida en el \autoref{sec:overloading}, considérense dos llamadas sobre el mismo circuito base.

\begin{lstlisting}[language={}, caption=Dos llamadas a la función sobrecargada at resueltas con distinta firma]
circuito1 = at 0 H base
circuito2 = at [0, 1, 2] H base
\end{lstlisting}

En la primera llamada, la generación de restricciones infiere que el literal \lstinline|0| tiene tipo \lstinline|i32|. Al consultar el vector de definiciones de \lstinline|at| en el ámbito, \lstinline|resolve| encuentra que \lstinline|i32| unifica con el primer parámetro de la firma \lstinline|i32 -> (qcirc -> qcirc) -> qcirc -> qcirc| y descarta la firma cuyo primer parámetro es \lstinline|[i32]|. En la segunda llamada ocurre lo contrario, ya que el literal de slice \lstinline|[0, 1, 2]| tiene tipo \lstinline|[i32]|, que solo unifica con la segunda firma. Ambas llamadas comparten el mismo nombre de función en el código fuente, pero el compilador selecciona una implementación distinta en cada punto de llamada sin ninguna anotación adicional por parte del programador.

\subsection{Detección de un error de tipos en tiempo de compilación}

El sistema de tipos rechaza en tiempo de compilación cualquier programa cuyas restricciones no puedan unificarse, incluso si la operación tendría sentido en un lenguaje de tipado dinámico. Considérese la siguiente declaración incorrecta.

\begin{lstlisting}[language={}, caption=Declaración que produce un error de unificación entre qbool e i32]
combinar :: qbool -> i32
combinar x = x + 1
\end{lstlisting}

La aplicación de \lstinline|+| exige, según su firma \lstinline|i32 -> i32 -> i32|, que su primer operando tenga tipo \lstinline|i32|, de modo que la generación de restricciones añade $\tau_x = \mathtt{i32}$. Sin embargo, la anotación de la declaración ya fijó $\tau_x = \mathtt{qbool}$ mediante la regla de inferencia de declaraciones. Al intentar unificar \lstinline|qbool| con \lstinline|i32| (\autoref{sec:unificacion}), ambos son tipos \lstinline|Builtin| con identificadores distintos, por lo que la unificación falla y el compilador emite un error señalando el span de \lstinline|x| dentro de la expresión \lstinline|x + 1|, junto con los dos tipos incompatibles. Este tipo de error, trivial de introducir por accidente en una biblioteca de tipado dinámico como Qiskit o Cirq (\autoref{sec:estadodelarte}), queda descartado en Qalc antes de que el programa llegue a generarse como circuito.

\chapter{Backend del compilador}
\label{ch:backend}

Este capítulo describe las tres fases finales del pipeline de Qalc, que transforman el \lstinline|TypedAST| ya verificado en el texto OpenQASM~3 que se ejecuta sobre hardware o simulador cuántico. El aplanador perezoso reduce el \lstinline|TypedAST| a un \lstinline|FlatAST|, una única expresión sin abstracción funcional pendiente. Sobre esa expresión aplanada, la comprobación de consumo lineal verifica que cada recurso cuántico se usa exactamente una vez, cerrando en tiempo de compilación la disciplina lineal presentada en el \autoref{sec:disciplina_lineal_uso}. Por último, la generación de código traduce el \lstinline|FlatAST| ya validado a instrucciones OpenQASM~3. La \autoref{tab:pipeline} sitúa estas tres fases dentro del pipeline completo del compilador.

\section{Aplanador perezoso}
\label{sec:inliner}

El aplanador perezoso transforma el \lstinline|TypedAST| en un \lstinline|FlatAST|, una representación aplanada orientada directamente a la generación de código. La motivación es que el código OpenQASM~3 que se emitirá en la fase siguiente no tiene noción de abstracción funcional, pues es una secuencia plana de instrucciones de puerta sobre registros concretos. El aplanador cierra esa brecha materializando, en tiempo de compilación, todas las aplicaciones de función.

El mecanismo consiste en interpretar el \lstinline|TypedAST| mediante una evaluación perezosa simbólica. El aplanador recorre el AST como si fuera un intérprete, pero en lugar de producir valores concretos, produce nodos del \lstinline|FlatAST|. El uso de thunks tiene un doble propósito. Por un lado, las expresiones que dependen de parámetros cuánticos de entrada, cuyo valor solo existe en tiempo de ejecución cuántica, se envuelven en thunks, es decir, nodos suspendidos del \lstinline|FlatAST| que representan un cómputo diferido y que el codegen traducirá directamente a secuencias de puertas. Por otro lado, la pereza permite omitir por completo el cómputo de las ramas o subexpresiones cuyo resultado nunca llega a demandarse, evitando así trabajo redundante que un aplanador estricto realizaría igualmente. Las expresiones cuyo valor puede determinarse completamente en tiempo de compilación (literales, aplicaciones de funciones conocidas sobre argumentos conocidos) se reducen directamente a su forma normal.

La estrategia perezosa es necesaria porque el lenguaje usa combinadores de orden superior como \lstinline|map| para expresar la aplicación repetida de una operación a un slice de cúbits. El aplanador despliega estos combinadores en tiempo de compilación conociendo la longitud del slice (que es estática, ya que el tamaño de los registros cuánticos es fijo), produciendo una secuencia plana de aplicaciones de puerta sin ninguna abstracción pendiente.

El resultado de esta fase es un único \lstinline|FlatAST|, una sola expresión que representa la totalidad del programa y que las dos fases siguientes de este capítulo recorren y transforman. Esa expresión contiene únicamente referencias a built-ins con sus argumentos concretos, composiciones secuenciales de operaciones, condicionales clásicos basados en resultados de medición y los thunks descritos a continuación. No existe ningún nodo lambda, ningún bloque \lstinline|let| ni ninguna referencia a funciones definidas por el usuario: todas han sido aplanadas y sus cuerpos sustituidos en el punto de llamada.

Un thunk es la unidad básica que sostiene la pereza del aplanador, es decir, una celda que representa un cómputo que puede estar todavía sin realizar. En Rust, un thunk se modela como una enumeración con un estado explícito, de modo que el aplanador puede distinguir en todo momento si una subexpresión ya está resuelta, si es una operación irreducible o si sigue pendiente de forzarse.

\begin{lstlisting}[language={}, caption=Esbozo del tipo Thunk que sostiene la pereza del aplanador]
pub enum Thunk {
    Value(FlatValue),
    Atomic(BuiltinCall),
    Pending(Box<Expression<Inferred>>),
    Executed(Box<FlatValue>),
}
\end{lstlisting}

Cada variante corresponde a uno de los estados que puede atravesar una subexpresión durante el aplanado.

\begin{itemize}
  \item \lstinline|Value(FlatValue)|. La subexpresión ya es un valor conocido en tiempo de compilación (un literal o el resultado de reducir una aplicación de built-ins sobre argumentos conocidos) y no requiere ninguna evaluación diferida.
  \item \lstinline|Atomic(BuiltinCall)|. La subexpresión es una única aplicación de built-in sobre argumentos ya resueltos, sin subexpresiones adicionales que evaluar; es la forma mínima que el codegen puede traducir directamente a una instrucción.
  \item \lstinline|Pending(Box<Expression<Inferred>>)|. La subexpresión todavía no se ha forzado. Guarda una referencia al fragmento del \lstinline|TypedAST| pendiente de interpretar, que el aplanador evaluará únicamente cuando otra parte del programa demande su resultado.
  \item \lstinline|Executed(Box<FlatValue>)|. El cómputo ya se forzó una vez y su resultado queda cacheado, de modo que una segunda demanda sobre el mismo thunk no repite el trabajo ni duplica las instrucciones emitidas.
\end{itemize}

La implementación actual de esta fase es bastante experimental. Cubre el aplanado de los combinadores de orden superior y el ciclo de vida del thunk descrito arriba para los casos más comunes del lenguaje, pero todavía no garantiza el aplanado correcto de todas las combinaciones posibles de bloques \lstinline|linear|, \lstinline|perform| y funciones de orden superior anidadas.

\section{Comprobación de consumo único de recursos lineales}
\label{sec:linear-checker}

El \autoref{sec:disciplina_lineal_uso} establece que toda variable cuántica debe consumirse exactamente una vez a lo largo de su ámbito, sin poder descartarse ni duplicarse. Esta fase es la que hace cumplir esa garantía en tiempo de compilación, recorriendo el \lstinline|FlatAST| producido por el aplanador antes de que la generación de código emita ninguna instrucción.

La comprobación se sitúa deliberadamente después del aplanador y no antes, sobre el \lstinline|TypedAST|. Antes de aplanar, un recurso lineal puede pasar de una función a otra a través de llamadas todavía no resueltas, de modo que contar sus usos exigiría seguir el flujo de valores a través de fronteras de función. Una vez aplanado el programa en una única expresión sin abstracciones pendientes, cada recurso lineal aparece directamente como argumento concreto de una aplicación de built-in, lo que permite contar sus apariciones de forma puramente sintáctica sobre el árbol.

El diseño de esta fase es análogo en espíritu al \autoref{sec:recursivity-checker}: es una validación pura que no produce una representación nueva, sino que aprueba o rechaza el \lstinline|FlatAST| tal cual, de modo que la transformación es \lstinline|FlatAST| $\to$ \lstinline|FlatAST|. El comprobador recorre la expresión única del \lstinline|FlatAST| y mantiene, para cada identificador de tipo lineal en el ámbito, un contador de apariciones. Un recurso que termina su ámbito con un contador en cero indica una fuga, es decir, un cúbit que nunca se midió ni se liberó. Un contador mayor que uno indica una duplicación, que equivaldría a clonar información cuántica, algo prohibido por el teorema de no clonación. En cualquiera de los dos casos se emite un error de compilación con el span del recurso implicado, indicando si el problema es una fuga o una duplicación.

Los thunks que introduce el aplanador reciben un tratamiento particular: un recurso lineal referenciado dentro de un thunk cuenta como usado en el punto en que el thunk se construye, no en el punto en que posteriormente se fuerza. De no ser así, un thunk forzado más de una vez internamente por el codegen produciría un conteo de usos incorrecto que no refleja el consumo real del programa fuente.

Al igual que el aplanador, la implementación actual de esta fase es bastante experimental. El conteo de apariciones descrito arriba cubre las expresiones que el aplanador reduce a los nodos habituales del \lstinline|FlatAST|, pero la interacción entre el conteo y los casos menos comunes de thunks anidados todavía no está completamente cubierta por pruebas.

\section{Generación de código}
\label{sec:codegen}

La fase de generación de código toma el \lstinline|FlatAST| producido por el aplanador y emite texto OpenQASM~3 válido. OpenQASM~3 es el estándar abierto de descripción de circuitos cuánticos, adoptado por IBM Quantum, IonQ y Rigetti, entre otros proveedores~\cite{openqasm3_2022}. A diferencia de versiones anteriores, OpenQASM~3 incorpora construcciones de control de flujo clásico (condicionales y bucles acotados basados en resultados de medición) que se corresponden directamente con los nodos del \lstinline|FlatAST|.

El diseño del codegen se estructura en dos partes. El preámbulo del fichero de salida declara la versión de OpenQASM~3, los registros cuánticos y clásicos necesarios, y las definiciones de las puertas derivadas de los built-ins que no están incluidas en el conjunto estándar de OpenQASM~3. Cada \lstinline|BuiltinDef| del crate \lstinline|qalc-builtins| tiene asociada una plantilla de emisión: para las operaciones aritméticas clásicas se emiten expresiones de OpenQASM~3; para las puertas cuánticas básicas (Hadamard, Pauli-X, CNOT, etc.) se emiten las instrucciones de puerta estándar del lenguaje.

El cuerpo del circuito se genera recorriendo el \lstinline|FlatAST| en orden topológico y emitiendo para cada nodo la instrucción OpenQASM~3 correspondiente: aplicaciones de puerta sobre cúbits concretos, mediciones con asignación a bits clásicos, y bloques \lstinline|if| condicionados a resultados de medición para el control de flujo clásico-cuántico. La correspondencia entre nodos del \lstinline|FlatAST| e instrucciones OpenQASM~3 es, por diseño, casi uno a uno; el trabajo de aplanar y especializar el programa ya ha sido realizado por el aplanador, de modo que el codegen puede limitarse a traducir sin necesidad de realizar transformaciones adicionales.

La emisión de cada built-in se resuelve mediante una extensión del tipo \lstinline|BuiltinDef| introducido en la \autoref{sec:organizacion-crates}, que además del nombre y la firma de tipo incorpora el propio mecanismo de traducción al formato de salida.

\begin{lstlisting}[language={}, caption=Esbozo de BuiltinDef parametrizado por backend y su emisión a OpenQASM 3]
pub trait Backend {}

pub struct OpenQasm3;
impl Backend for OpenQasm3 {}

pub struct Quil;
impl Backend for Quil {}

pub trait Transpilable<B: Backend> {
    fn transpile(&self, args: &[FlatValue]) -> String;
}

pub struct BuiltinDef<B: Backend> {
    pub name: &'static str,
    pub signature: TypeWithUnknow,
    pub body: Box<dyn Transpilable<B>>,
}

pub struct OpenQasm3Transpilable {
    pub emit: fn(&[FlatValue]) -> String,
}

impl Transpilable<OpenQasm3> for OpenQasm3Transpilable {
    fn transpile(&self, args: &[FlatValue]) -> String {
        (self.emit)(args)
    }
}
\end{lstlisting}

\lstinline|Backend| es un rasgo marcador, sin métodos propios, que identifica un formato de salida soportado por el compilador. \lstinline|OpenQasm3| es el struct que lo implementa para el backend actual; \lstinline|Quil|, el lenguaje de ensamblado cuántico de Rigetti, se incluye como ejemplo de un segundo backend hipotético para ilustrar la extensión. El rasgo \lstinline|Transpilable| se parametriza sobre ese \lstinline|Backend|, de modo que un mismo built-in puede tener una implementación de transpilación distinta para cada formato de salida sin que ambas compartan tipo. \lstinline|BuiltinDef| se parametriza igualmente, así que el compilador podría mantener un \lstinline|BuiltinDef<OpenQasm3>| y, en el futuro, un \lstinline|BuiltinDef<Quil>| para el mismo built-in, cada uno con su propia lógica de emisión, sin duplicar el resto de la definición. Todo built-in, sin excepción, se define mediante una implementación concreta del rasgo \lstinline|Transpilable|, cuyo único método, \lstinline|transpile|, recibe los argumentos ya aplanados del built-in y devuelve directamente el fragmento de código que lo representa. \lstinline|OpenQasm3Transpilable| es la implementación concreta que usa el backend actual e implementa \lstinline|Transpilable<OpenQasm3>|, envolviendo simplemente la función Rust que compone el fragmento de OpenQASM~3 de salida para ese built-in en particular, ya sea el nombre de una instrucción de puerta estándar o una expresión aritmética compuesta.

Parametrizar \lstinline|Transpilable| y \lstinline|BuiltinDef| sobre el backend responde directamente al objetivo de diseño, ya mencionado en el \autoref{ch:arquitectura}, de mantener la arquitectura del compilador modular y extensible. Añadir soporte para un backend nuevo no exigiría modificar \lstinline|BuiltinDef| ni \lstinline|Transpilable|, ni tocar \lstinline|qalc-builtins| ni el resto del pipeline, sino únicamente definir el struct que marca el nuevo backend y aportar, para cada built-in, la implementación de \lstinline|Transpilable| correspondiente a ese backend.

El cuerpo de un built-in es código Rust ordinario y no una expresión de Qalc, por lo que nunca pasa por las fases de comprobación de tipos ni de consumo lineal descritas en los capítulos anteriores, que solo recorren el AST del programa del usuario. Esta asimetría es deliberada, ya que al quedar el cuerpo de los built-ins fuera del alcance de esas validaciones, la implementación de un built-in puede incorporar comportamientos ocultos o convenientes que no serían expresables ni permitidos en código Qalc ordinario, como gestionar registros clásicos o cúbits auxiliares que el programador no declaró explícitamente, o consultar contexto adicional del proceso de emisión. Esta libertad es la que permitirá, en implementaciones futuras del rasgo \lstinline|Transpilable|, resolver built-ins que necesiten comportamientos especiales para funcionar correctamente sin comprometer la disciplina que sí se exige al resto del programa.

\chapter{Conclusiones y líneas futuras}
\label{ch:conclusiones}

\section{Conclusiones}

Este trabajo ha abordado el diseño de un lenguaje de programación cuántica de alto nivel y la implementación de su compilador, partiendo de una observación central expuesta en el \autoref{sec:contexto}, que ningún lenguaje existente combina un sistema de tipos lineales estático con una salida hacia un formato estándar e interoperable con hardware de múltiples fabricantes. Los dos objetivos de diseño que se derivaron de esa observación han guiado el desarrollo completo de Qalc.

El primer objetivo, dotar al lenguaje de un sistema de tipos lineal sin la penalización de rendimiento que introducen las soluciones embebidas en un lenguaje anfitrión como Guppy, se ha alcanzado combinando la inferencia de tipos de Hindley-Milner con la disciplina de consumo único descrita en el \autoref{ch:analisis-semantico}. Al implementarse en un compilador de propósito específico y no como una biblioteca embebida en un lenguaje de tipado dinámico, Qalc verifica en tiempo de compilación tanto la clonación accidental como la fuga de cúbits sin liberar, sin necesidad de anotaciones de tipo explícitas en la mayoría de las expresiones del programa. La sobrecarga de funciones discriminada por tipo, descrita en el \autoref{sec:overloading}, complementa este sistema de tipos permitiendo expresar operaciones polimórficas sencillas sin recurrir a clases de tipos.

El segundo objetivo, generar código hacia un formato estándar e interoperable, se ha alcanzado con la fase de generación de código descrita en el \autoref{ch:backend}, que traduce el programa aplanado a OpenQASM~3. El diseño del mecanismo de emisión, basado en el rasgo \lstinline|Transpilable| parametrizado por el tipo de backend, se planteó explícitamente para que añadir un formato de salida adicional no exigiera modificar el resto del compilador, en línea con el objetivo de modularidad y extensibilidad que atraviesa toda la arquitectura descrita en el \autoref{ch:arquitectura}.

El resultado conjunto de ambos objetivos es un prototipo funcional en el que el frontend del compilador, desde el análisis léxico hasta la inferencia de tipos y la verificación de recursividad, está completamente implementado. Las fases del backend (el aplanador perezoso, la comprobación de consumo lineal y la generación de código, descritas en el \autoref{ch:backend}) se encuentran en un estado experimental, suficiente para demostrar el funcionamiento completo del pipeline sobre programas que respetan la disciplina lineal descrita en el \autoref{sec:disciplina_lineal_uso}, pero todavía no cubren la totalidad de las combinaciones posibles del lenguaje.

Tal y como se anticipó en el \autoref{sec:alcance}, este trabajo no aborda la descomputación automática de estructuras de datos cuánticos recursivos. Se trata de una limitación deliberada y no de un objetivo incumplido, ya que ese problema exige un análisis estático del grafo topológico del estado cuántico que excede el alcance de un prototipo de estas características y que constituye, como ya se argumentó, una línea de investigación propia.

\section{Líneas futuras}

El diseño modular del compilador, en el que cada fase se comunica con la siguiente mediante el rasgo \lstinline|CompilerStep| descrito en el \autoref{ch:arquitectura}, permite identificar líneas de trabajo futuro que se pueden abordar de forma incremental sin rediseñar el compilador existente.

\begin{itemize}
  \item \textit{Descomputación automática de estructuras recursivas.} Extender el análisis semántico para sintetizar automáticamente el operador de descomputación de estructuras de datos cuánticos recursivos, como se discutió en el \autoref{sec:alcance}, combinando análisis de alcance estático con una representación enriquecida del grafo de entrelazamiento.

  \item \textit{Maduración de las fases del backend.} Completar la cobertura del aplanador perezoso sobre combinaciones anidadas de combinadores de orden superior y bloques monádicos, y extender las pruebas de la comprobación de consumo lineal a los casos de thunks anidados descritos en el \autoref{sec:linear-checker}, hasta alcanzar el mismo nivel de madurez que el frontend del compilador.

  \item \textit{Backends adicionales.} Aprovechar la separación entre el rasgo \lstinline|Transpilable| y su implementación concreta para OpenQASM~3, descrita en el \autoref{sec:codegen}, para añadir soporte a formatos de salida adicionales, como Quil, validando en la práctica la extensibilidad que motivó ese diseño.

  \item \textit{Sobrecarga de funciones más expresiva.} Generalizar la sobrecarga discriminada por el tipo del primer parámetro, descrita en el \autoref{sec:overloading}, hacia mecanismos de resolución que consideren varios parámetros o que se aproximen a las clases de tipos, siempre que la inferencia de tipos siga siendo decidible.

  \item \textit{Evaluación empírica sobre hardware cuántico.} Ejecutar los circuitos generados por Qalc sobre simuladores y hardware físico real de distintos proveedores compatibles con OpenQASM~3, para medir la profundidad y la fidelidad de los circuitos resultantes frente a implementaciones equivalentes en Qiskit o Cirq.
\end{itemize}

\chapter{Summary and Conclusions}
\label{ch:conclusions}

\section{Conclusions}

This work has addressed the design of a high-level quantum programming language and the implementation of its compiler, starting from a central observation laid out in \autoref{sec:contexto}, that no existing language combines a static linear type system with output to a standard, interoperable format supported by multiple hardware vendors. The two design objectives that followed from that observation have guided the entire development of Qalc.

The first objective, giving the language a linear type system without the performance penalty introduced by solutions embedded in a host language such as Guppy, has been met by combining Hindley-Milner type inference with the single-consumption discipline described in \autoref{ch:analisis-semantico}. By being implemented as a purpose-built compiler rather than a library embedded in a dynamically typed language, Qalc verifies at compile time both accidental cloning and unreleased qubit leaks, without requiring explicit type annotations on most expressions of the program. Type-discriminated function overloading, described in \autoref{sec:overloading}, complements this type system by allowing simple polymorphic operations to be expressed without resorting to type classes.

The second objective, generating code towards a standard, interoperable format, has been met with the code generation phase described in \autoref{ch:backend}, which translates the flattened program into OpenQASM~3. The design of the emission mechanism, based on the \lstinline|Transpilable| trait parameterized over the backend type, was deliberately conceived so that adding a further output format would not require modifying the rest of the compiler, in line with the modularity and extensibility goal that runs through the whole architecture described in \autoref{ch:arquitectura}.

The combined result of both objectives is a functional prototype in which the compiler frontend, from lexical analysis through type inference and recursion verification, is fully implemented. The backend phases (the lazy flattener, the linear consumption check and the code generation, described in \autoref{ch:backend}) remain at an experimental stage, sufficient to demonstrate the complete operation of the pipeline over programs that respect the linear discipline described in \autoref{sec:disciplina_lineal_uso}, but they do not yet cover every possible combination of the language.

As anticipated in \autoref{sec:alcance}, this work does not address automatic uncomputation. This is a deliberate limitation rather than an unmet objective, since that problem requires a static analysis of the topological graph of the quantum state that exceeds the scope of a prototype of this kind and constitutes, as already argued, a research line of its own.

\section{Future work}

The modular design of the compiler, in which each phase communicates with the next through the \lstinline|CompilerStep| trait described in \autoref{ch:arquitectura}, makes it possible to identify future lines of work that can be tackled incrementally without redesigning the existing compiler.

\begin{itemize}
  \item \textit{Automatic uncomputation.} Extend the semantic analysis to automatically synthesize the uncomputation operator for recursive quantum data structures, as discussed in \autoref{sec:alcance}, combining static scope analysis with a richer representation of the entanglement graph.

  \item \textit{Maturing the backend phases.} Complete the coverage of the lazy flattener over nested combinations of higher-order combinators and monadic blocks, and extend the test coverage of the linear consumption check to the nested-thunk cases described in \autoref{sec:linear-checker}, until it reaches the same level of maturity as the compiler frontend.

  \item \textit{Expanding the standard library of types and operations.} The delivered prototype offers an illustrative catalog of built-ins (the basic quantum gates, a small set of arithmetic operations and the functions used in this document's examples), sufficient to demonstrate the language design but not for productive use. Completing this standard library relies on the same \lstinline|BuiltinDef| and \lstinline|Transpilable| mechanism already described in \autoref{ch:backend}, without requiring structural changes to the compiler.

  \item \textit{Additional backends.} Exploit the separation between the \lstinline|Transpilable| trait and its concrete implementation for OpenQASM~3, described in \autoref{sec:codegen}, to add support for further output formats, such as Quil, validating in practice the extensibility that motivated that design.

  \item \textit{More expressive function overloading.} Generalize the overloading discriminated by the type of the first parameter, described in \autoref{sec:overloading}, towards resolution mechanisms that consider several parameters or that approach type classes, as long as type inference remains decidable.

  \item \textit{Empirical evaluation on quantum hardware.} Run the circuits generated by Qalc on simulators and on real physical hardware from different providers compatible with OpenQASM~3, to measure the depth and fidelity of the resulting circuits against equivalent implementations in Qiskit or Cirq.
\end{itemize}

\chapter{Presupuesto}
\label{ch:presupuesto}

Este capítulo recoge una estimación del coste total del proyecto, desglosada en recursos materiales y coste de mano de obra. Todos los importes se expresan en euros. La suscripción a herramientas de inteligencia artificial se contrató en dólares estadounidenses y se ha convertido aplicando un tipo de cambio aproximado de 1~USD~$\approx$~0,92~EUR.

\section{Recursos y herramientas}

El equipamiento principal del proyecto es el ordenador portátil empleado para el desarrollo. Su coste se amortiza proporcionalmente a la duración del proyecto, estimada en seis meses, sobre una vida útil de cuatro años (48 meses), lo que resulta en una fracción amortizable de 6/48. Adicionalmente, durante tres meses del desarrollo se utilizó una suscripción mensual a herramientas de asistencia mediante inteligencia artificial, contratada a 20~\$/mes (aproximadamente 18,50~€/mes al tipo de cambio indicado).

\begin{table}[!ht]
\centering
\caption{Recursos y herramientas}
\label{tbl:recursos}
\begin{tabularx}{0.85\textwidth}{Xcr}
    \toprule
    \textbf{Descripción} & \textbf{Cantidad} & \textbf{Coste (€)} \\
    \midrule
    Portátil de desarrollo (amortización 6/48 meses sobre 1.000,00~€) & 1 & 125,00 \\
    Suscripción a herramientas de IA (3 meses × 18,50~€/mes) & 3 meses & 55,50 \\
    \midrule
    \textbf{Subtotal} & & \textbf{180,50} \\
    \bottomrule
\end{tabularx}
\end{table}

\section{Mano de obra}

El desarrollo del proyecto ha requerido aproximadamente 300 horas de trabajo distribuidas en tres fases: investigación del dominio y diseño del lenguaje, implementación del compilador en Rust y redacción de la presente memoria. La tarifa aplicada es de 15,00~€ por hora, correspondiente a un perfil de ingeniero de software junior en España.

\begin{table}[!ht]
\centering
\caption{Coste de mano de obra}
\label{tbl:mano-de-obra}
\begin{tabularx}{0.85\textwidth}{Xccr}
    \toprule
    \textbf{Descripción} & \textbf{Horas} & \textbf{Tarifa (€/h)} & \textbf{Coste (€)} \\
    \midrule
    Investigación y diseño & 80 & 15,00 & 1.200,00 \\
    Desarrollo del compilador & 150 & 15,00 & 2.250,00 \\
    Redacción de la memoria & 70 & 15,00 & 1.050,00 \\
    \midrule
    \textbf{Total} & \textbf{300} & & \textbf{4.500,00} \\
    \bottomrule
\end{tabularx}
\end{table}

\section{Coste total}

\begin{table}[!ht]
\centering
\caption{Resumen del presupuesto del proyecto}
\label{tbl:presupuesto}
\begin{tabularx}{0.85\textwidth}{Xr}
    \toprule
    \textbf{Concepto} & \textbf{Coste (€)} \\
    \midrule
    Recursos y herramientas & 180,50 \\
    Mano de obra & 4.500,00 \\
    \midrule
    \textbf{Coste total del proyecto} & \textbf{4.680,50} \\
    \bottomrule
\end{tabularx}
\end{table}


\appendix

\chapter{Gramática del lenguaje en notación EBNF}
\label{ch:gramatica-ebnf}

Este apéndice recoge la especificación sintáctica completa del lenguaje Qalc. La gramática se transcribe fielmente desde el fichero \lstinline|grammar.pest|, que es la fuente de verdad que emplea el analizador sintáctico (véase el \autoref{sec:parser}), y se expresa en notación EBNF en su variante ISO/Wirth.

\section{Convenciones de notación}
\label{sec:ebnf-convenciones}

Cada regla tiene la forma \lstinline|nombre = expresión ;|, donde el signo igual asocia un nombre a su definición y el punto y coma cierra la regla. Las alternativas se separan con \lstinline{|}, la concatenación se indica con una coma, la agrupación con paréntesis, la repetición de cero o más veces con llaves \lstinline|{ }|, y la parte opcional con corchetes \lstinline|[ ]|. Los terminales del lenguaje se escriben entre comillas dobles, de modo que los caracteres \lstinline|{|, \lstinline|}|, \lstinline|[| y \lstinline|]| que forman parte de la sintaxis de Qalc aparecen entrecomillados y no se confunden con la metasintaxis. La excepción se denota con el signo menos, es decir, \lstinline|A - B| describe lo que encaja con \lstinline|A| pero no con \lstinline|B|. Las clases de caracteres que no pueden desglosarse cómodamente se describen entre interrogantes siguiendo la convención de secuencia especial de ISO, y las aclaraciones intercaladas se escriben entre \lstinline|(*| y \lstinline|*)|.

La transcripción conserva los nombres de regla originales, pero conviene tener presentes algunas particularidades del formalismo de partida que la notación EBNF no captura de manera directa. Pest implementa gramáticas de expresión de análisis, por lo que la alternativa se resuelve de forma ordenada, probando cada opción de izquierda a derecha y quedándose con la primera que encaja, mientras que en EBNF la alternativa es conceptualmente no ordenada. Las reglas léxicas (atómicas en Pest) no admiten espacios en su interior, en cambio entre los símbolos de las reglas sintácticas se permite de forma implícita cualquier \lstinline|whitespace| o \lstinline|comment|. La búsqueda negativa de Pest, que no forma parte de EBNF, se aproxima mediante la excepción con el signo menos o con una aclaración. Por último, los anclajes \lstinline|SOI| y \lstinline|EOI| marcan respectivamente el inicio y el final de la entrada.

A lo largo de la gramática se emplean los siguientes auxiliares para las clases de caracteres.

\begin{lstlisting}[language=EBNF, caption=Clases de caracteres auxiliares]
letter          = ? carácter alfabético ASCII, a-z o A-Z ? ;
digit           = ? dígito decimal ASCII, 0-9 ? ;
hex_digit       = ? dígito hexadecimal ASCII, 0-9 a-f A-F ? ;
newline         = ? salto de línea ? ;
tab             = ? tabulador ? ;
carriage_return = ? retorno de carro ? ;
any_character   = ? cualquier carácter ? ;
end_of_input    = ? fin de la entrada (EOI) ? ;
\end{lstlisting}

\section{Espacios en blanco y comentarios}
\label{sec:ebnf-espacios}

\begin{lstlisting}[language=EBNF, caption=Reglas de espacios en blanco y comentarios]
whitespace    = " " | tab | carriage_return | newline ;
block_comment = "/*" , { block_comment | ( any_character - "*/" ) } , "*/" ;
line_comment  = ( "//" | "#" ) , { any_character - newline } ,
                ( newline | end_of_input ) ;
comment       = line_comment | block_comment ;
\end{lstlisting}

\section{Palabras reservadas e identificadores}
\label{sec:ebnf-identificadores}

\begin{lstlisting}[language=EBNF, caption=Reglas de palabras reservadas e identificadores]
keyword           = "if" | "else" | "let" | "type" | "in" | "then"
                  | "match" | "perform" | "linear" | "where" | "return" ;
simple_identifier = ( ( letter | "_" ) , { letter | digit | "_" } )
                  - keyword ;
\end{lstlisting}

\section{Literales}
\label{sec:ebnf-literales}

\begin{lstlisting}[language=EBNF, caption=Reglas de literales numéricos]
dec_integer    = digit , { digit } ;
hex_integer    = "0x" , hex_digit , { hex_digit } ;
bin_integer    = "0b" , ( "0" | "1" ) , { "0" | "1" } ;
unknow_integer = hex_integer | bin_integer | dec_integer ;
u8_integer     = ( hex_integer | bin_integer | dec_integer ) , "u8" ;
u16_integer    = ( hex_integer | bin_integer | dec_integer ) , "u16" ;
u32_integer    = ( hex_integer | bin_integer | dec_integer ) , "u32" ;
u64_integer    = ( hex_integer | bin_integer | dec_integer ) , "u64" ;
i8_integer     = ( hex_integer | bin_integer | dec_integer ) , "i8" ;
i16_integer    = ( hex_integer | bin_integer | dec_integer ) , "i16" ;
i32_integer    = ( hex_integer | bin_integer | dec_integer ) , "i32" ;
i64_integer    = ( hex_integer | bin_integer | dec_integer ) , "i64" ;
float_exponent = ( "e" | "E" ) , [ "+" | "-" ] , digit , { digit } ;
float_decimal  = digit , { digit } , "." , digit , { digit } ;
float_leading  = "." , digit , { digit } ;
float_number   = float_decimal , [ float_exponent ]
               | float_leading , [ float_exponent ]
               | ( digit , { digit } ) , float_exponent ;
unknow_float   = float_number ;
f32_float      = ( float_number | ( digit , { digit } ) ) , "f32" ;
f64_float      = ( float_number | ( digit , { digit } ) ) , "f64" ;
literal        = f32_float | f64_float | unknow_float
               | u8_integer | u16_integer | u32_integer | u64_integer
               | i8_integer | i16_integer | i32_integer | i64_integer
               | unknow_integer ;
\end{lstlisting}

\section{Operadores}
\label{sec:ebnf-operadores}

\begin{lstlisting}[language=EBNF, caption=Reglas de operadores]
apply_operator         = "$" ;
pipe_operator          = "|>" ;
logic_or_operator      = "||" ;
logic_and_operator     = "&&" ;
or_operator            = "|" ;
and_operator           = "&" ;
xor_operator           = "^" ;
equal_operator         = "==" ;
not_equal_operator     = "!=" ;
less_equal_operator    = "<=" ;
greater_equal_operator = ">=" ;
less_operator          = "<" ;
greater_operator       = ">" ;
add_operator           = "+" ;
subtract_operator      = "-" ;
multiply_operator      = "*" ;
divide_operator        = "/" ;
module_operator        = "%" ;
raised_operator        = "**" ;
overloadable_operators = logic_or_operator | logic_and_operator
               | equal_operator | not_equal_operator
               | less_equal_operator | greater_equal_operator
               | less_operator | greater_operator
               | add_operator | subtract_operator | raised_operator
               | multiply_operator | divide_operator | module_operator
               | or_operator | and_operator | xor_operator ;
binary_operators = pipe_operator | apply_operator
               | logic_or_operator | logic_and_operator
               | equal_operator | not_equal_operator
               | less_equal_operator | greater_equal_operator
               | less_operator | greater_operator
               | add_operator | subtract_operator | raised_operator
               | multiply_operator | divide_operator | module_operator
               | or_operator | and_operator | xor_operator ;
\end{lstlisting}

\section{Expresiones átomo}
\label{sec:ebnf-atomos}

\begin{lstlisting}[language=EBNF, caption=Reglas de las expresiones átomo]
operator_identifier     = "`" , overloadable_operators , "`" ;
identifier              = simple_identifier | operator_identifier ;
pattern                 = simple_identifier
                        | "(" , pattern ,
                          ( "," , pattern ) , { "," , pattern } , ")" ;
partial_operation_left  = "(" , binary_operators , expression_atom , ")" ;
partial_operation_right = "(" , expression_atom , binary_operators , ")" ;
call_expression         = ( identifier | ( "(" , expression , ")" ) ) ,
                          call_operand , { call_operand } ;
    (* call_operand agrupa el "+" de Pest; no es una regla del .pest *)
call_operand            = block_expression | lambda_expression
                        | partial_operation_right | partial_operation_left
                        | ( "(" , expression , ")" ) | tuple_expression
                        | slice_expression | identifier | literal ;
lambda_expression       = "(" , { pattern } , ")" , "->" , expression ;
tuple_expression        = ( "(" , ")" )
                        | "(" , expression ,
                          ( "," , expression ) , { "," , expression } , ")" ;
slice_range_expression  = "[" , expression , ".." , expression , "]" ;
slice_items_expression  = "[" , [ expression , { "," , expression } ] , "]" ;
slice_expression        = slice_items_expression | slice_range_expression ;
\end{lstlisting}

\section{Expresiones de bloque}
\label{sec:ebnf-bloques}

\begin{lstlisting}[language=EBNF, caption=Reglas de las expresiones de bloque]
assigment_expression = pattern , "=" , expression ;
let_block            = "let" , assigment_expression ,
                       { "," , assigment_expression } , "in" , expression ;
if_else_block        = "if" , expression , "then" , expression ,
                       "else" , expression ;
match_block          = "match" , expression , "{" ,
                       ( expression , "->" , expression ) ,
                       { "," , expression , "->" , expression } , "}" ;
return_expression    = "return" , expression ;
monadic_assigment    = pattern , "<-" , expression ;
monadic_block        = "{" ,
                       ( assigment_expression | monadic_assigment | expression ) ,
                       { "," ,
                         ( assigment_expression | monadic_assigment | expression ) } ,
                       [ "," , return_expression ] , "}" ;
perform_block        = "perform" , monadic_block ;
linear_block         = "linear" , monadic_block ;
block_expression     = let_block | if_else_block | match_block
                     | perform_block | linear_block ;
\end{lstlisting}

\section{Expresiones binarias}
\label{sec:ebnf-binarias}

\begin{lstlisting}[language=EBNF, caption=Reglas de las expresiones binarias]
expression_atom = block_expression | lambda_expression
               | partial_operation_right | partial_operation_left
               | call_expression | ( "(" , expression , ")" )
               | tuple_expression | slice_expression | identifier | literal ;
expression      = expression_atom , { binary_operators , expression_atom } ;
\end{lstlisting}

\section{Tipos}
\label{sec:ebnf-tipos}

\begin{lstlisting}[language=EBNF, caption=Reglas del sistema de tipos]
type_inner       = "(" , type , ")" ;
type_slice       = "[" , type , "]" ;
type_tuple       = ( "(" , type ,
                     ( "," , type ) , { "," , type } , ")" )
                 | ( "(" , ")" ) ;
type_with_params = simple_identifier , "<" , type_atom , { type_atom } , ">" ;
type_atom        = type_with_params | simple_identifier | type_slice
                 | type_tuple | type_inner ;
type             = type_atom , { "->" , type_atom } ;
\end{lstlisting}

\section{Declaraciones y producción inicial}
\label{sec:ebnf-declaraciones}

\begin{lstlisting}[language=EBNF, caption=Reglas de declaraciones y producción inicial]
type_function_annotation = identifier , "::" , type ;
type_restriction         = "linear" ;
type_declaration         = "type" , simple_identifier , "{" ,
                           type_function_annotation ,
                           { type_function_annotation } , "}" ,
                           [ "where" , type_restriction ,
                             { "," , type_restriction } ] ,
                           [ ";" ] ;
function_declaration     = type_function_annotation , identifier ,
                           { pattern } , "=" , expression , ";" ;
declarations             = type_declaration | function_declaration ;
program                  = { declarations } ;
    (* delimitado por SOI (inicio) y EOI (fin de la entrada) *)
\end{lstlisting}


%%%%%%%%%%%%%%%%%%%%%%%%%%%%%%%%%%%%%%%%%%%%%%%%%%%%%%%%%%
% Bibliografía
%%%%%%%%%%%%%%%%%%%%%%%%%%%%%%%%%%%%%%%%%%%%%%%%%%%%%%%%%%
\backmatter

\printbibliography[prenote=note]
%%%%%%%%%%%%%%%%%%%%%%%%%%%%%%%%%%%%%%%%%%%%%%%%%%%%%%%%%%

\end{document}

